% NOTE: This file is included in the document body (not the preamble).
% Load required packages in the main preamble instead, e.g.:
%   \usepackage{xcolor,multicol,amssymb,amsmath}
% (expl3 is part of the LaTeX kernel in modern LaTeX.)

%  ANSWERS TO EXERCISES
%  The Physics Codex: A Complete University Foundation
%  Korede Samuel Omotosho (Falcon)
%  Aurikrex Learning Systems, 2026
%
%  File: back_matter/answers.tex
% ============================================================

\chapter*{Answers to Exercises}
\addcontentsline{toc}{chapter}{Answers to Exercises}
\markboth{Answers to Exercises}{Answers to Exercises}

\small
% Answers contain long numerical expressions; permit modest interword
% adjustment so prose wraps naturally without hiding genuine equation overflow.
\setlength{\emergencystretch}{3em}

% -------------------- Formatting helpers (answers only) --------------------
\makeatletter
\newif\ifAnswers@FirstChapter
\Answers@FirstChaptertrue
\newif\ifAnswers@FirstB
\Answers@FirstBtrue

\definecolor{navyblue}{RGB}{26, 39, 68}

% Fallbacks (in case the main preamble doesn't load amsmath/amssymb).
\providecommand{\text}[1]{\mathrm{#1}}
\providecommand{\checkmark}{\ensuremath{\surd}}
\providecommand{\vect}[1]{\mathbf{#1}}

\newcommand{\AnswersChapterBand}[2]{%
  \ifAnswers@FirstChapter\global\Answers@FirstChapterfalse\else\vspace{14pt}\fi
  \noindent
  \colorbox{navyblue}{%
    \parbox{\linewidth}{%
      \vspace{6pt}%
      \hspace{8pt}{\color{white}\large\bfseries Chapter #1 \quad #2}%
      \vspace{6pt}%
    }%
  }%
  \vspace{10pt}%
  \global\Answers@FirstBtrue
}

\newcommand{\AnswersSectionHeading}[1]{%
  \vspace{8pt}%
  \noindent{\bfseries #1}%
  \vspace{4pt}%
}

\let\Answers@origtextbf\textbf

% Special \textbf handler for A- and B- labels (keeps normal \textbf{} elsewhere).
\ExplSyntaxOn
\cs_new_protected:Npn \Answers_textbf:n #1
  {
    % A-labels like "A1.~C" (the "~" is just a space token here, so match whitespace)
    \regex_match:nnTF { \A\s*A\d+\.\s*[A-D]\s*\Z } {#1}
      {
        \regex_extract_once:nnN { \A\s*A(\d+)\.\s*([A-D])\s*\Z } {#1} \l_tmpa_seq
        % Print label as "A1." then option letter; DO NOT insert \quad here
        % because the document already has "\textbf{A1.~C} \quad ...".
        \noindent\Answers@origtextbf{A\seq_item:Nn \l_tmpa_seq {1}.}~\seq_item:Nn \l_tmpa_seq {2}
      }
      {
        % B-labels like "B1."
        \regex_match:nnTF { \A\s*B\d+\.\s*\Z } {#1}
          {
            \regex_extract_once:nnN { \A\s*B(\d+)\.\s*\Z } {#1} \l_tmpb_seq
            \ifAnswers@FirstB\global\Answers@FirstBfalse\else
              \medskip
              \noindent\rule{\linewidth}{0.3pt}
              \medskip
            \fi
            \vspace{6pt}
            \noindent\Answers@origtextbf{B\seq_item:Nn \l_tmpb_seq {1}.}\par
            \vspace{3pt}
          }
          {\Answers@origtextbf{#1}}
      }
  }

% Alias without expl3 “:” so it can be used with \let in normal LaTeX code.
\cs_new_eq:NN \AnswersTextbf \Answers_textbf:n
\ExplSyntaxOff

% (a), (b), ... part labels in Section B
\newcommand{\Answers@partlabel}[1]{\noindent\Answers@origtextbf{(#1)}\space}
\begingroup\lccode`\~=`(\lowercase{\endgroup\def~}{\Answers@paren}
\ExplSyntaxOn
\cs_set_protected:Npn \Answers@paren #1)
  {
    \mode_if_math:TF
      { (#1) }
      {
        % Only treat (a), (b), ... as part labels; everything else is a normal parenthesis.
        \regex_match:nnTF { \A\s*[A-Za-z]\s*\Z } {#1}
          {
            \Answers@partlabel{#1}%
            \@ifnextchar~{\@gobble}{}%
          }
          { (#1) }
      }
  }
\ExplSyntaxOff
\makeatother

% Public alias (no @) so it can be used outside \makeatletter, e.g. when making '(' active.
\makeatletter
\let\AnswersParen\Answers@paren
\makeatother

% -------------------- Intro paragraph --------------------
\begin{center}
\textit{Section A answers are given as the correct option letter followed by a brief justification.
Section B answers give the key numerical result or the essential reasoning for each part; full worked
solutions for selected problems are available on the Aurikrex Learning Platform.}
\end{center}
\vspace{12pt}

% ============================================================
%  CHAPTER 1
% ============================================================
\AnswersChapterBand{1}{Measurement}

\AnswersSectionHeading{Section A}
\begin{multicols}{2}
\begingroup
\let\textbf\AnswersTextbf
\def\quad{ --- }
\setlength{\parindent}{0pt}

\textbf{A1.~C} \quad Electric current (ampere) is one of
the seven SI base quantities. Force, velocity, and energy
are all derived.

\textbf{A2.~B} \quad
$3.5\times10^{-4}\,\text{m} = 3.5\times10^{-4}\times10^{3}
\,\text{mm} = 0.35\,\text{mm}$.

\textbf{A3.~A} \quad
$P = F/A \Rightarrow [MLT^{-2}]/[L^{2}] = [ML^{-1}T^{-2}]$.

\textbf{A4.~B} \quad $0.004\,070$: digits $4,0,7,0$ are
significant (the leading zeros are not). That gives
\textbf{4} significant figures.

\textbf{A5.~B} \quad $108\,\text{km\,h}^{-1}
\div 3.6 = 30\,\text{m\,s}^{-1}$.

\textbf{A6.~C} \quad Momentum is mass$\times$velocity,
a vector. Temperature, power and mass are scalars.

\textbf{A7.~B} \quad
$[m][g][h]=[M][LT^{-2}][L]=[ML^{2}T^{-2}]$, the
dimension of energy. \checkmark

\textbf{A8.~B} \quad Five nearly identical readings
$(24.3,\,24.3,\,24.4,\,24.3,\,24.3)$ are
\textit{precise} (small spread). Without a true value
for comparison we cannot confirm accuracy.

\textbf{A9.~C} \quad
$0.85\,\text{g\,cm}^{-3}\times1000 = 850\,\text{kg\,m}^{-3}$.

\textbf{A10.~C} \quad $E=mva$ has dimensions
$[M][LT^{-1}][LT^{-2}]=[ML^{2}T^{-3}]$ (power), not
energy $[ML^{2}T^{-2}]$. The equation is dimensionally
inconsistent.

\endgroup
\end{multicols}

\AnswersSectionHeading{Section B}
\begingroup
\let\textbf\AnswersTextbf
\setlength{\parindent}{0pt}
\def\quad{\par\smallskip}
% (Parenthesis activation removed: it broke ordinary parentheses in text/math.)

\textbf{B1.}
(a)~$[F/A]=[MLT^{-2}]/[L^{2}]=[ML^{-1}T^{-2}]$. \quad
(b)~$[W/t]=[ML^{2}T^{-2}]/[T]=[ML^{2}T^{-3}]$.
Watts: $\text{kg\,m}^{2}\text{s}^{-3}$. \quad
(c)~$n=[V/(kT)]=[ML^{2}T^{-2}]/([ML^{2}T^{-2}K^{-1}][K])
=[M^{0}L^{0}T^{0}]$, dimensionless. Consistent. \quad
(d)~$k=PV/(nT)$; units
$\text{Pa}\cdot\text{m}^{3}/(\text{mol}\cdot\text{K})
=\text{J\,mol}^{-1}\text{K}^{-1}$.

\textbf{B2.}
(a)~Absolute error $\pm0.05\,\text{cm}$; percentage
error $=(0.05/12.4)\times100=0.40\%$. \quad
(b)~$\rho=m/V=210/(25.0\times10^{-6})
=8.40\times10^{6}\,\text{g\,m}^{-3}
=8.40\times10^{3}\,\text{kg\,m}^{-3}$. \quad
(c)~Fractional uncertainty in $V=3\times\delta L/L
=3\times0.5/25=6\%$; in $m$ negligible;
$\delta\rho/\rho=6\%$,
$\delta\rho=0.06\times8400=504
\approx500\,\text{kg\,m}^{-3}$. \quad
(d)~Report as $\rho=(8400\pm500)\,\text{kg\,m}^{-3}$.

\textbf{B3.}
(a)~$T=2\pi\sqrt{L/g}\Rightarrow
g=4\pi^{2}L/T^{2}$; dimensions $[L]/[T^{2}]=[LT^{-2}]$.
\checkmark \quad
(b)~$g=4\pi^{2}(0.650)/(1.618)^{2}
=4\times9.870\times0.650/2.618
=9.79\,\text{m\,s}^{-2}$. \quad
(c)~$\%\,\text{err in }g
=\%\,\text{err in }L + 2\times\%\,\text{err in }T
=(0.5/65.0)\times100+2\times(0.5/161.8)\times100
=0.77+0.62=1.39\approx1.4\%$. \quad
(d)~$\delta g=0.014\times9.79=0.14\,\text{m\,s}^{-2}$.
Report: $g=(9.79\pm0.14)\,\text{m\,s}^{-2}$.

\textbf{B4.}
(a)~Mean $=\frac{1}{5}(5.23+5.25+5.21+5.24+5.22)
=26.15/5=5.23\,\text{cm}$. \quad
(b)~Deviations: $0,+0.02,-0.02,+0.01,-0.01$;
mean absolute deviation $=0.06/5=0.012\,\text{cm}$.
\quad
(c)~Percentage uncertainty $=(0.012/5.23)\times100
=0.23\%$. \quad
(d)~Systematic error shifts all readings by the same
amount (e.g.\ a zero error on the ruler); random errors
produce spread in both directions. Random errors are
reduced by repeating and averaging; systematic errors
are not.

\textbf{B5.}
(a)~$[k]=[F/x]=[MLT^{-2}]/[L]=[MT^{-2}]$. \quad
(b)~$[T]=[(m/k)^{1/2}]=[M/(MT^{-2})]^{1/2}=[T^{2}]^{1/2}
=[T]$. Consistent. \quad
(c)~$k=4\pi^{2}m/T^{2}$. \quad
(d)~Plot $T^{2}$ vs $m$; gradient $=4\pi^{2}/k$,
hence $k=4\pi^{2}/\text{gradient}$. \quad
(e)~$k=4\times9.870/0.452^{2}
\approx39.48/0.2043=193\,\text{N\,m}^{-1}$.

% ============================================================
%  CHAPTER 2
% ============================================================
\endgroup

\AnswersChapterBand{2}{Scalars and Vectors}

\AnswersSectionHeading{Section A}
\begin{multicols}{2}
\begingroup
\let\textbf\AnswersTextbf
\def\quad{ --- }
\setlength{\parindent}{0pt}

\textbf{A1.~B} \quad Displacement and velocity are both
vectors (magnitude + direction). Speed is scalar; mass
is scalar; power is scalar.

\textbf{A2.~D} \quad Resultant range:
$|8-6|=2\,\text{N}$ to $8+6=14\,\text{N}$.
$15\,\text{N}$ lies outside this range and is impossible.

\textbf{A3.~A} \quad $F_x = 100\cos60°=100\times0.5
=50\,\text{N}$.

\textbf{A4.~C} \quad $5^{2}+12^{2}=25+144=169=13^{2}$.
This is a Pythagorean triple, so the angle between the
$5\,\text{N}$ and $12\,\text{N}$ forces is $90°$.

\textbf{A5.~B} \quad
$|\vect{F}|=\sqrt{(-3)^{2}+(-4)^{2}}=\sqrt{25}=5\,\text{N}$.
Magnitude is always positive.

\textbf{A6.~C} \quad $v_{P\text{ rel }Q}=v_P-v_Q
=40\,\text{N}-(-40\,\text{N})=80\,\text{km\,h}^{-1}$
north.

\textbf{A7.~A} \quad Resultant $R=\sqrt{F_1^2+F_2^2
+2F_1F_2\cos\theta}$; maximum when $\cos\theta=1$,
i.e.\ $\theta=0°$.

\textbf{A8.~B} \quad
$|\vect{s}|=\sqrt{5^2+(-12)^2}=\sqrt{169}=13\,\text{m}$.

\textbf{A9.~B} \quad Two equal forces $F$ at $90°$:
resultant $=\sqrt{F^2+F^2}=F\sqrt{2}$.

\textbf{A10.~B} \quad By the triangle of forces, when
three concurrent forces are in equilibrium the resultant
of any two must equal the third in magnitude and act
in the opposite direction.

\endgroup
\end{multicols}

\AnswersSectionHeading{Section B}
\begingroup
\let\textbf\AnswersTextbf
\setlength{\parindent}{0pt}
\def\quad{\par\smallskip}
% (Parenthesis activation removed: it broke ordinary parentheses in text/math.)

\textbf{B1.}
(a)~$F_x=80\cos30°=69.3\,\text{N}$;
$F_y=80\sin30°=40\,\text{N}$. \quad
(b)~$R_x=69.3+60=129.3\,\text{N}$;
$R_y=40-45=-5\,\text{N}$. \quad
(c)~$|\vect{R}|=\sqrt{129.3^2+5^2}=\sqrt{16718.5}
=129.4\,\text{N}\approx129\,\text{N}$. \quad
(d)~$\theta=\tan^{-1}(5/129.3)=2.2°$ below the
horizontal.

\textbf{B2.}
(a)~By Lami's theorem, for three concurrent coplanar
forces in equilibrium:
$T_1/\sin\angle_1=T_2/\sin\angle_2=W/\sin\angle_3$. \quad
(b)~The angle between $T_1$ and $T_2$ (above the
horizontal) and $W$ (downward):
$T_1=W\sin60°/\sin90°=200\times0.866/1=173.2\,\text{N}$;
$T_2=W\sin30°/\sin90°=200\times0.5=100\,\text{N}$. \quad
(c)~Both tensions reduce as the angle between the strings
decreases; in the limit (strings horizontal) tension
would become infinite.

\textbf{B3.}
(a)~$R_{horizontal}=40\cos0°+30\cos(-90°)+50\cos120°
=40+0-25=15\,\text{N}$ east. \quad
$R_{vertical}=40\sin0°+30\sin(-90°)+50\sin120°
=0-30+43.3=13.3\,\text{N}$ north. \quad
(b)~$|\vect{R}|=\sqrt{15^2+13.3^2}=\sqrt{402}
=20.1\,\text{N}$. \quad
(c)~$\theta=\tan^{-1}(13.3/15)=41.6°$ north of east.

\textbf{B4.}
(a)~$v_{AB}=v_A-v_B$. $v_A=50\,\text{km\,h}^{-1}$
north; $v_B=40\,\text{km\,h}^{-1}$ east. \quad
$|v_{AB}|=\sqrt{50^2+40^2}=\sqrt{4100}=64\,\text{km\,h}^{-1}$. \quad
(b)~Direction: $\theta=\tan^{-1}(40/50)=38.7°$ west of
north. \quad
(c)~If $v_{river}=3\,\text{m\,s}^{-1}$ east and
$v_{boat in still water}=5\,\text{m\,s}^{-1}$, to travel
straight across: aim upstream at angle
$\theta=\sin^{-1}(3/5)=36.9°$ west of north; resultant
speed $=\sqrt{25-9}=4\,\text{m\,s}^{-1}$.

\textbf{B5.}
(a)~$T\cos\theta=mg$ and $T\sin\theta=F$. Hence
$\tan\theta=F/(mg)$. \quad
(b)~$T=\sqrt{F^2+(mg)^2}$. \quad
(c)~$\theta=\tan^{-1}[0.50/(0.200\times10)]
=\tan^{-1}(0.25)=14.0°$. \quad
(d)~$T=\sqrt{0.50^2+2.0^2}=\sqrt{4.25}=2.06\,\text{N}$. \quad
(e)~If a second horizontal force $F_2=0.30\,\text{N}$
acts perpendicular to $F_1$: resultant horizontal force
$=\sqrt{0.50^2+0.30^2}=0.583\,\text{N}$;
$\theta=\tan^{-1}(0.583/2.0)=16.3°$;
$T=\sqrt{0.583^2+2.0^2}=2.08\,\text{N}$.

% ============================================================
%  CHAPTER 3
% ============================================================
\endgroup

\AnswersChapterBand{3}{Kinematics}

\AnswersSectionHeading{Section A}
\begin{multicols}{2}
\begingroup
\let\textbf\AnswersTextbf
\def\quad{ --- }
\setlength{\parindent}{0pt}

\textbf{A1.~B} \quad $s=\tfrac{1}{2}at^2\Rightarrow
100=\tfrac{1}{2}\times a\times25\Rightarrow a=8\,\text{m\,s}^{-2}$.

\textbf{A2.~C} \quad Horizontal line on a $v$-$t$ graph
means $v$ is constant, i.e.\ zero acceleration.

\textbf{A3.~C} \quad At height $H=u^2/2g$: speed is zero.
At half-height $H/2$: $v^2=u^2-2g(H/2)=u^2-u^2/2=u^2/2$,
so $v=u/\sqrt{2}$. Ratio of speed at $H/2$ to speed at
$H$ is $(u/\sqrt{2}):0=\sqrt{2}:0$.

\textbf{A4.~B} \quad Range $R=v^2\sin2\theta/g$.
At $45°$: $R_{45}=v^2/g$.
At $30°$: $R_{30}=v^2\sin60°/g=R_{45}\times(\sqrt{3}/2)
\approx0.87R$.

\textbf{A5.~B} \quad $H=\tfrac{1}{2}gT^2$.
For $4H$: $T'^2=4T^2\Rightarrow T'=2T$.

\textbf{A6.~C} \quad Area under a $v$-$t$ graph equals
displacement (a vector, not merely distance).

\textbf{A7.~C} \quad In projectile motion (no air
resistance) horizontal and vertical motions are
completely independent.

\textbf{A8.~C} \quad Average speed in first $4\,\text{s}$:
$15\,\text{m\,s}^{-1}$. Next $4\,\text{s}$:
$25\,\text{m\,s}^{-1}$. Speed increasing $\Rightarrow$
accelerating.

\textbf{A9.~B} \quad $s=\tfrac{1}{2}(u+v)t
=\tfrac{1}{2}(30+0)(6)=90\,\text{m}$.

\textbf{A10.~C} \quad Velocity (m\,s$^{-1}$) and speed
(m\,s$^{-1}$) share the same SI unit.

\endgroup
\end{multicols}

\AnswersSectionHeading{Section B}
\begingroup
\let\textbf\AnswersTextbf
\setlength{\parindent}{0pt}
\def\quad{\par\smallskip}
% (Parenthesis activation removed: it broke ordinary parentheses in text/math.)

\textbf{B1.}
(a)~Using $v^2=u^2+2as$: $20^2=0+2a(400)\Rightarrow
a=\frac{400}{800}=0.5\,\text{m\,s}^{-2}$. \quad
(b)~Braking: $0^2=20^2+2(-3)s\Rightarrow s=200/3
=66.7\,\text{m}$. \quad
(c)~Reaction distance $=20\times0.8=16\,\text{m}$.
Total stopping distance $=16+66.7=82.7\,\text{m}$. \quad
(d)~At $v=25\,\text{m\,s}^{-1}$: reaction distance
$=25\times0.8=20\,\text{m}$; braking distance
$=25^2/(2\times3)=104.2\,\text{m}$; total
$=124.2\,\text{m}$.

\textbf{B2.}
(a)~$u_x=v_0\cos40°=25\times0.766=19.15\,\text{m\,s}^{-1}$;
$u_y=25\sin40°=16.07\,\text{m\,s}^{-1}$. \quad
(b)~$t_{top}=u_y/g=16.07/10=1.607\,\text{s}$. \quad
(c)~$H=u_y^2/(2g)=258.2/20=12.91\,\text{m}$. \quad
(d)~$t_{flight}=2t_{top}=3.21\,\text{s}$. \quad
(e)~$R=u_x\times t_{flight}=19.15\times3.21=61.5\,\text{m}$.

\textbf{B3.}
(a)~From rest: area under $v$-$t$ graph (triangles/trapezoids)
for each phase. Phase 1 (0--10s): triangle, area
$=\tfrac{1}{2}\times10\times20=100\,\text{m}$.
Phase 2 (10--25s): rectangle, $15\times20=300\,\text{m}$.
Phase 3 (25--35s): triangle, $\tfrac{1}{2}\times10\times20
=100\,\text{m}$. Total $=500\,\text{m}$. \quad
(b)~Acceleration in phase 1: $a=20/10=2\,\text{m\,s}^{-2}$.
Deceleration in phase 3: $a=-20/10=-2\,\text{m\,s}^{-2}$.
\quad
(c)~Average speed $=500/35=14.3\,\text{m\,s}^{-1}$.

\textbf{B4.}
(a)~Time to fall $h=80\,\text{m}$:
$h=\tfrac{1}{2}gt^2\Rightarrow t=\sqrt{160/10}=4\,\text{s}$. \quad
(b)~Horizontal range $=u\times t=15\times4=60\,\text{m}$. \quad
(c)~Vertical speed at impact: $v_y=gt=10\times4=40\,\text{m\,s}^{-1}$.
Resultant speed $=\sqrt{15^2+40^2}=\sqrt{1825}=42.7\,\text{m\,s}^{-1}$. \quad
(d)~Angle to horizontal $=\tan^{-1}(40/15)=69.4°$.

\textbf{B5.}
(a)~Let the trains each have length $L$ and speed $v$.
Relative speed when approaching $=2v$.
Time to pass $=2L/(2v)=L/v$.
\quad
(b)~Relative speed when same direction $=v-v=0$;
the trains take infinitely long to pass (they travel
together). For one stationary and one moving at $v$:
time $=2L/v$. \quad
(c)~$v_{rel}=30+20=50\,\text{m\,s}^{-1}$ (head-on).
Time to pass $=(250+200)/50=450/50=9\,\text{s}$. \quad
(d)~Same direction: $v_{rel}=30-20=10\,\text{m\,s}^{-1}$.
Time $=450/10=45\,\text{s}$.

% ============================================================
%  CHAPTER 4
% ============================================================
\endgroup

\AnswersChapterBand{4}{Dynamics}

\AnswersSectionHeading{Section A}
\begin{multicols}{2}
\begingroup
\let\textbf\AnswersTextbf
\def\quad{ --- }
\setlength{\parindent}{0pt}

\textbf{A1.~C} \quad Impulse $=Ft=20\times5=100\,\text{N\,s}
=\Delta p$.

\textbf{A2.~B} \quad $a=(m_2-m_1)g/(m_1+m_2)
=(6-2)\times10/(6+2)=40/8=5\,\text{m\,s}^{-2}$.

\textbf{A3.~B} \quad Newton's Third Law pair: the book
pulls Earth upward with the same force that Earth pulls
the book downward. Weight and normal force are
\textit{not} Newton 3rd-Law pairs (they act on the
same object).

\textbf{A4.~B} \quad $a=g\sin37°=10\times0.6
=6\,\text{m\,s}^{-2}$ (smooth incline).

\textbf{A5.~C} \quad Terminal velocity when
$\text{drag}=\text{weight}$; net force $=0$;
acceleration $=0$.

\textbf{A6.~B} \quad Perfectly inelastic collision:
objects stick together. Momentum is conserved; kinetic
energy is not (some becomes heat/sound/deformation).

\textbf{A7.~B} \quad Maximum static friction
$=\mu N=0.4\times100=40\,\text{N}$. The applied force
$30\,\text{N}<40\,\text{N}$, so the block does not slide.
Friction equals applied force $=30\,\text{N}$.

\textbf{A8.~C} \quad Apparent weight
$=m(g+a)=60(10+2)=720\,\text{N}$.

\textbf{A9.~C} \quad Conservation of momentum:
$m_{bullet}v_{bullet}=m_{rifle}v_{rifle}$;
$v_{rifle}=0.015\times600/3.0=3.0\,\text{m\,s}^{-1}$.

\textbf{A10.~C} \quad Passengers lurching forward when
the bus brakes: their inertia resists the deceleration
--- Newton's First Law.

\endgroup
\end{multicols}

\AnswersSectionHeading{Section B}
\begingroup
\let\textbf\AnswersTextbf
\setlength{\parindent}{0pt}
\def\quad{\par\smallskip}
% (Parenthesis activation removed: it broke ordinary parentheses in text/math.)

\textbf{B1.}
(a)~Net force $=800-200=600\,\text{N}$;
$a=F/m=600/1200=0.5\,\text{m\,s}^{-2}$. \quad
(b)~$v=u+at=0+0.5\times15=7.5\,\text{m\,s}^{-1}$. \quad
(c)~$s=\tfrac{1}{2}at^2=\tfrac{1}{2}\times0.5\times225
=56.25\,\text{m}$. \quad
(d)~Normal force from road: $N=mg=1200\times10
=12\,000\,\text{N}$; friction provides centripetal force.

\textbf{B2.}
(a)~Both blocks: net force $=(5-3)\times10=20\,\text{N}$;
total mass $=8\,\text{kg}$; $a=20/8=2.5\,\text{m\,s}^{-2}$. \quad
(b)~Tension in string: for $3\,\text{kg}$ block:
$T-30=3\times2.5\Rightarrow T=37.5\,\text{N}$. \quad
(c)~String goes slack if $m_2\leq m_1$; tension then
zero.

\textbf{B3.}
(a)~Momentum before: $0.2\times300=60\,\text{N\,s}$.
After: $(0.2+5)\times v=60\Rightarrow v=60/5.2
=11.5\,\text{m\,s}^{-1}$. \quad
(b)~$\Delta KE=\tfrac{1}{2}\times5.2\times11.5^2
-\tfrac{1}{2}\times0.2\times300^2=344-9000=-8656\,\text{J}$.
Lost as heat and deformation. \quad
(c)~Coefficient of friction: $\mu=v^2/(2gs)=11.5^2/(2\times10\times s)$.
If $s=2\,\text{m}$: $\mu=132.25/40=3.3$ (unrealistically
high, showing the block does not travel far).

\textbf{B4.}
(a)~On the $5\,\text{kg}$ block:
$mg\sin30°-f=ma\Rightarrow5\times10\times0.5-f=5a
\Rightarrow25-f=5a$.
Normal: $N=mg\cos30°=5\times10\times0.866=43.3\,\text{N}$.
Friction: $f=\mu N=0.3\times43.3=13.0\,\text{N}$.
$a=(25-13.0)/5=2.4\,\text{m\,s}^{-2}$. \quad
(b)~$v^2=2\times2.4\times3.0=14.4\Rightarrow v=3.79\,\text{m\,s}^{-1}$.
\quad
(c)~At the bottom the block leaves the ramp; projectile
motion follows.

\textbf{B5.}
(a)~$F_{net}=F_{thrust}-F_{weight}-F_{drag}
=8000-3000-500=4500\,\text{N}$.
$a=4500/300=15\,\text{m\,s}^{-2}$. \quad
(b)~To hover: thrust $=$ weight $=3000\,\text{N}$. \quad
(c)~As fuel burns, mass decreases; same thrust produces
greater acceleration. The rocket accelerates faster
as its mass decreases. \quad
(d)~$v_{exhaust}\,\Delta m=M\,\Delta v$ (rocket equation);
momentum of ejected gas equals momentum gained by rocket.

% ============================================================
%  CHAPTER 5
% ============================================================
\endgroup

\AnswersChapterBand{5}{Work, Energy and Power}

\AnswersSectionHeading{Section A}
\begin{multicols}{2}
\begingroup
\let\textbf\AnswersTextbf
\def\quad{ --- }
\setlength{\parindent}{0pt}

\textbf{A1.~A} \quad $W=Fd\cos\theta
=50\times8.0\times\cos60°=400\times0.5=200\,\text{J}$.

\textbf{A2.~C} \quad Kinetic energy $=\tfrac{1}{2}mv^2$
(by definition of KE).

\textbf{A3.~B} \quad $E=\tfrac{1}{2}kx^2
=\tfrac{1}{2}\times400\times(0.05)^2=0.50\,\text{J}$.

\textbf{A4.~B} \quad $P=\dot{m}gh=50\times10\times10
=5000\,\text{W}$.

\textbf{A5.~C} \quad Power of engine $=$ engine force
$\times$ velocity $=3000\times20=60\,000\,\text{W}$.

\textbf{A6.~C} \quad Efficiency $=P_{out}/P_{in}
=1.5/2=0.75=75\%$.

\textbf{A7.~D} \quad Gravity acts toward the Sun
(centripetally); velocity is tangential (perpendicular to
gravity). Work done $=Fd\cos90°=0$.

\textbf{A8.~B} \quad Energy conservation:
$\tfrac{1}{2}mv^2=\tfrac{1}{2}mv_0^2+mgH\Rightarrow
v=\sqrt{v_0^2+2gH}$.

\textbf{A9.~A} \quad $KE=\tfrac{1}{2}mv^2$.
Doubling $KE$: $v_{new}=v\sqrt{2}$, so speed is
multiplied by $\sqrt{2}$.

\textbf{A10.~B} \quad $P=mgh/t=70\times10\times12/18
=8400/18=466.7\approx467\,\text{W}$.
\endgroup
\end{multicols}

\AnswersSectionHeading{Section B}
\begingroup
\let\textbf\AnswersTextbf
\setlength{\parindent}{0pt}
\def\quad{\par\smallskip}
% (Parenthesis activation removed: it broke ordinary parentheses in text/math.)


\textbf{B1.}
(a)~Normal force $N=mg\cos20°=5\times10\times0.940=47.0\,\text{N}$.
Friction $=\mu N=0.3\times47.0=14.1\,\text{N}$.
Net force $=mg\sin20°-f=5\times10\times0.342-14.1
=17.1-14.1=3.0\,\text{N}$.
$a=0.60\,\text{m\,s}^{-2}$. \quad
(b)~Work by gravity over $s=4\,\text{m}$:
$W_g=mgs\sin20°=5\times10\times4\times0.342=68.4\,\text{J}$.
Work by friction: $W_f=-14.1\times4=-56.4\,\text{J}$.
Net work $=68.4-56.4=12.0\,\text{J}=\Delta KE$. \checkmark \quad
(c)~$v=\sqrt{2\times0.60\times4}=2.19\,\text{m\,s}^{-1}$.

\textbf{B2.}
(a)~Spring compressed $x=F/k=200/5000=0.04\,\text{m}$;
$E_{spring}=\tfrac{1}{2}\times5000\times0.04^2=4\,\text{J}$. \quad
(b)~At launch: $4=\tfrac{1}{2}\times0.05\times v^2
\Rightarrow v=\sqrt{160}=12.65\,\text{m\,s}^{-1}$. \quad
(c)~Maximum height: $h=v^2/(2g)=160/20=8\,\text{m}$. \quad
(d)~At $h=3\,\text{m}$: $KE=4-mgh=4-0.05\times10\times3
=4-1.5=2.5\,\text{J}$.

\textbf{B3.}
(a)~$P_{useful}=Fv=1200\times25=30\,000\,\text{W}$.
Engine power $=30\,000/0.35=85\,714\approx85.7\,\text{kW}$. \quad
(b)~At constant speed, engine force $=$ resistance force
$=1200\,\text{N}$. \quad
(c)~Work done against resistance per hour
$=1200\times25\times3600=1.08\times10^8\,\text{J}$.

\textbf{B4.}
(a)~$PE_{top}=mgh=0.5\times10\times3=15\,\text{J}$;
$KE_{top}=0$; total $E=15\,\text{J}$. \quad
(b)~At bottom: $KE=15\,\text{J}$;
$v=\sqrt{2\times15/0.5}=\sqrt{60}=7.75\,\text{m\,s}^{-1}$. \quad
(c)~With friction $W_{friction}=3\,\text{J}$:
$KE_{bottom}=15-3=12\,\text{J}$;
$v=\sqrt{2\times12/0.5}=6.93\,\text{m\,s}^{-1}$. \quad
(d)~Efficiency $=12/15=80\%$.

\textbf{B5.}
(a)~Work done per hour $=Power\times time
=75\times3600=270\,000\,\text{J}$. \quad
(b)~Stairs: $P=mgh/t=75\times10\times5/20
=187.5\,\text{W}$; cycling $=75\,\text{W}$ assumed
mechanical output. Stairs requires $2.5\times$ the
power. \quad
(c)~Calories: $1\,\text{kcal}=4186\,\text{J}$.
In $1\,\text{h}$ at $75\,\text{W}$:
$270\,000/4186=64.5\,\text{kcal}$.

% ============================================================
%  CHAPTER 6
% ============================================================
\endgroup

\AnswersChapterBand{6}{Circular Motion}

\AnswersSectionHeading{Section A}
\begin{multicols}{2}
\begingroup
\let\textbf\AnswersTextbf
\def\quad{ --- }
\setlength{\parindent}{0pt}

\textbf{A1.~B} \quad Speed is constant; velocity changes
direction (hence changes). No net work done, but a
centripetal force acts.

\textbf{A2.~C} \quad $F=mv^2/r=0.2\times16/0.5
=6.4\,\text{N}$.

\textbf{A3.~C} \quad Kepler: $T^2\propto r^3$.
$T_2^2/T_1^2=(4r)^3/r^3=64\Rightarrow T_2=8T$.

\textbf{A4.~B} \quad $v_{max}=\sqrt{\mu rg}
=\sqrt{0.4\times50\times10}=\sqrt{200}=14.1\,\text{m\,s}^{-1}$.

\textbf{A5.~B} \quad At top of circle, minimum speed
when tension $T=0$: $mg=mv^2/r\Rightarrow v=\sqrt{rg}$.

\textbf{A6.~C} \quad $\theta=\omega t=10\times5
=50\,\text{rad}$.

\textbf{A7.~C} \quad At the ideal banking angle for a
given speed, the horizontal component of normal force
alone provides centripetal force. Friction is zero.

\textbf{A8.~C} \quad $v=\sqrt{ar}=\sqrt{0.0027
\times3.84\times10^8}=\sqrt{1.037\times10^6}
\approx1018\approx1020\,\text{m\,s}^{-1}$.

\textbf{A9.~B} \quad At hilltop: $mg-N=mv^2/r$.
Maximum speed when $N=0$: $v=\sqrt{rg}
=\sqrt{40\times10}=20\,\text{m\,s}^{-1}$.

\textbf{A10.~B} \quad At top: $(T+mg)=mv^2/L$.
If $T=mg$: $2mg=mv^2/L\Rightarrow v=\sqrt{2gL}$.

\endgroup
\end{multicols}

\AnswersSectionHeading{Section B}
\begingroup
\let\textbf\AnswersTextbf
\setlength{\parindent}{0pt}
\def\quad{\par\smallskip}
% (Parenthesis activation removed: it broke ordinary parentheses in text/math.)

\textbf{B1.}
(a)~$\omega=2\pi f=2\pi\times3=18.85\,\text{rad\,s}^{-1}$. \quad
(b)~$v=\omega r=18.85\times0.5=9.42\,\text{m\,s}^{-1}$. \quad
(c)~$a_c=\omega^2 r=18.85^2\times0.5=177.7\,\text{m\,s}^{-2}$. \quad
(d)~$F=ma_c=0.3\times177.7=53.3\,\text{N}$.

\textbf{B2.}
(a)~$T\cos\theta=mg$ and $T\sin\theta=m\omega^2 r$
where $r=L\sin\theta$. Dividing:
$\tan\theta=\omega^2L\sin\theta/g\Rightarrow
\cos\theta=g/(\omega^2 L)$. \quad
(b)~$\omega=\sqrt{g/(L\cos\theta)}$. \quad
(c)~$L=0.50\,\text{m}$, $\theta=30°$:
$\omega=\sqrt{10/(0.50\times\cos30°)}
=\sqrt{10/0.433}=\sqrt{23.1}=4.80\,\text{rad\,s}^{-1}$;
$n=\omega/2\pi=0.764\,\text{rev\,s}^{-1}=45.8\,\text{rpm}$.

\textbf{B3.}
(a)~Centripetal force: $F=mv^2/r=1200\times400/80
=6000\,\text{N}$. \quad
(b)~Friction available $=\mu mg=0.6\times1200\times10
=7200\,\text{N}>6000\,\text{N}$. Safe. \quad
(c)~Max speed: $v=\sqrt{\mu rg}=\sqrt{0.6\times80\times10}
=\sqrt{480}=21.9\,\text{m\,s}^{-1}=78.9\,\text{km\,h}^{-1}$.

\textbf{B4.}
(a)~$T^2=4\pi^2R^3/(GM_E)$;
$M_E=4\pi^2R^3/(GT^2)
=4\times9.87\times(6.4\times10^6+400\times10^3)^3/
(6.67\times10^{-11}\times(5535)^2)$. \quad
(b)~Orbital speed:
$v=2\pi r/T=2\pi\times6.78\times10^6/5535
=7700\,\text{m\,s}^{-1}=7.7\,\text{km\,s}^{-1}$. \quad
(c)~Period $\propto r^{3/2}$; higher orbit $\Rightarrow$
longer period. GEO period $=24\,\text{h}$.

\textbf{B5.}
(a)~At bottom: $N-mg=mv^2/r\Rightarrow N=m(g+v^2/r)
=70\times(10+400/20)=70\times30=2100\,\text{N}$. \quad
(b)~At top: $N+mg=mv^2/r\Rightarrow N=m(v^2/r-g)
=70\times(20-10)=700\,\text{N}$. \quad
(c)~Normal force at bottom is $3.5\times$ body weight;
at top is $1\times$ body weight --- rollercoaster riders
feel heaviest at the bottom.

% ============================================================
%  CHAPTER 7
% ============================================================
\endgroup

\AnswersChapterBand{7}{Simple Harmonic Motion}

\AnswersSectionHeading{Section A}
\begin{multicols}{2}
\begingroup
\let\textbf\AnswersTextbf
\def\quad{ --- }
\setlength{\parindent}{0pt}

\textbf{A1.~C} \quad SHM: $a\propto-x$
(acceleration proportional to displacement,
directed toward equilibrium).

\textbf{A2.~C} \quad $T=2\pi\sqrt{L/g}$.
If $g'=4g$: $T'=2\pi\sqrt{L/(4g)}=T/2$.

\textbf{A3.~B} \quad $v=\omega\sqrt{A^2-x^2}$;
maximum when $x=0$ (at equilibrium).

\textbf{A4.~B} \quad $T'=2\pi\sqrt{m'/k'}
=2\pi\sqrt{2m/(k/2)}=2\pi\sqrt{4m/k}=2T$.

\textbf{A5.~C} \quad $a_{max}=\omega^2A=25\times0.04
=1.00\,\text{m\,s}^{-2}$.

\textbf{A6.~C} \quad KE $=$ PE when
$\tfrac{1}{2}m\omega^2(A^2-x^2)=\tfrac{1}{2}m\omega^2x^2
\Rightarrow x=A/\sqrt{2}$.

\textbf{A7.~B} \quad Critical damping: system returns
to equilibrium as fast as possible without oscillating.

\textbf{A8.~B} \quad Resonance: driving frequency equals
natural frequency of the system.

\textbf{A9.~B} \quad $T=2\pi\sqrt{L/g}
=2\pi\sqrt{1.0/10}=2\pi/\sqrt{10}$.
Using $\pi^2=10$: $T=2\sqrt{10}/\sqrt{10}=2.0\,\text{s}$.

\textbf{A10.~C} \quad $x=A\sin\omega t$ (displacement);
$a=-\omega^2A\sin\omega t=-\omega^2 x$.
Phase difference between $x$ and $a$ is $180°$.
\endgroup
\end{multicols}

\AnswersSectionHeading{Section B}
\begingroup
\let\textbf\AnswersTextbf
\setlength{\parindent}{0pt}
\def\quad{\par\smallskip}
% (Parenthesis activation removed: it broke ordinary parentheses in text/math.)


\textbf{B1.}
(a)~$A=3.0\,\text{cm}=0.03\,\text{m}$;
$\omega=2\pi/T=2\pi/0.8=7.854\,\text{rad\,s}^{-1}$. \quad
(b)~$v_{max}=\omega A=7.854\times0.03=0.236\,\text{m\,s}^{-1}$. \quad
(c)~$a_{max}=\omega^2A=61.7\times0.03=1.85\,\text{m\,s}^{-2}$. \quad
(d)~$x=0.03\sin(7.854t)$ m; $v=0.236\cos(7.854t)$ m\,s$^{-1}$. \quad
(e)~At $x=0.02\,\text{m}$:
$v=\omega\sqrt{A^2-x^2}=7.854\sqrt{0.0009-0.0004}
=7.854\times0.02236=0.1756\,\text{m\,s}^{-1}$.

\textbf{B2.}
(a)~$T=2\pi\sqrt{L/g}$; $L=g(T/2\pi)^2
=10\times(2/2\pi)^2=10/(pi^2)=1.013\,\text{m}$. \quad
(b)~On Mars ($g=3.72\,\text{m\,s}^{-2}$):
$T=2\pi\sqrt{1.013/3.72}=3.28\,\text{s}$. \quad
(c)~The pendulum gains time on Earth (oscillates faster)
and loses time on Mars (slower). \quad
(d)~For a clock: change in $g$ by $\Delta g$ causes
fractional change in rate $\approx\Delta g/(2g)$.

\textbf{B3.}
(a)~$k=m\omega^2=0.5\times(2\pi/0.6)^2=0.5\times109.7
=54.8\,\text{N\,m}^{-1}$. \quad
(b)~$v_{max}=\omega A=10.47\times0.08=0.838\,\text{m\,s}^{-1}$. \quad
(c)~$E_{total}=\tfrac{1}{2}kA^2=\tfrac{1}{2}\times54.8
\times0.0064=0.175\,\text{J}$. \quad
(d)~At $x=A/2=0.04\,\text{m}$:
$KE=E-\tfrac{1}{2}k(0.04)^2=0.175-0.0438=0.131\,\text{J}$
(75\% of total).

\textbf{B4.}
(a)~$T=2\pi/\omega=2\pi/12=0.524\,\text{s}$;
$f=1/T=1.91\,\text{Hz}$. \quad
(b)~$a=-\omega^2 x=-144\times0.05=-7.2\,\text{m\,s}^{-2}$
(toward equilibrium). \quad
(c)~$v=\omega\sqrt{A^2-x^2}=12\sqrt{0.25-0.0025}
=12\times0.4975=5.97\,\text{m\,s}^{-1}$.

\textbf{B5.}
(a)~Forced oscillations: external driver supplies energy.
If driver frequency $\neq$ natural frequency, oscillation
amplitude is limited. \quad
(b)~At resonance: amplitude becomes very large (limited
only by damping). \quad
(c)~Tacoma Narrows Bridge (1940): wind-induced forced
oscillation matched natural frequency, causing
catastrophic resonance and collapse. \quad
(d)~Damping reduces resonant amplitude, widens resonance
peak, shifts resonant frequency slightly lower.
Lightly damped: tall, sharp peak; heavily damped: low,
broad peak.

% ============================================================
%  CHAPTER 8
% ============================================================
\endgroup

\AnswersChapterBand{8}{Gravitation}

\AnswersSectionHeading{Section A}
\begin{multicols}{2}
\begingroup
\let\textbf\AnswersTextbf
\def\quad{ --- }
\setlength{\parindent}{0pt}

\textbf{A1.~D} \quad $F\propto1/r^2$; doubling $r$
gives $F/4$.

\textbf{A2.~C} \quad $g=GM/R^2$.
$g'=G(2M)/(2R)^2=2GM/4R^2=g/2$.

\textbf{A3.~C} \quad $v_{esc}=\sqrt{2GM/R}$.
Same $M$, double $R$: $v'=v_{esc}/\sqrt{2}
=11.2/\sqrt{2}=7.92\,\text{km\,s}^{-1}$.

\textbf{A4.~C} \quad $T^2\propto r^3$.
$T_Y^2/T_X^2=(4r_X)^3/r_X^3=64$;
$T_Y=8\times8=64$ years.

\textbf{A5.~C} \quad Both astronaut and ISS are in free
fall (both falling toward Earth at the same rate);
apparent gravity is zero.

\textbf{A6.~B} \quad Gravitational potential = work done
per unit mass bringing a test mass from infinity to
that point.

\textbf{A7.~C} \quad $v=\sqrt{GM/r}$; if $r$
(approximately) doubles for large $h$, $v$
decreases by factor $\sqrt{2}$.

\textbf{A8.~B} \quad Kepler's Second Law: equal areas in
equal times; fastest when closest (perihelion).

\textbf{A9.~B} \quad
$V=-GM_E/R_E=-gR_E=-9.81\times6.4\times10^6
=-6.28\times10^7\,\text{J\,kg}^{-1}$.
For $1\,\text{kg}$: $PE=-6.28\times10^7\,\text{J}$.

\textbf{A10.~B} \quad Geostationary: above equator,
period $24\,\text{h}$, orbits west to east. Appears
stationary from Earth.

\endgroup
\end{multicols}

\AnswersSectionHeading{Section B}
\begingroup
\let\textbf\AnswersTextbf
\setlength{\parindent}{0pt}
\def\quad{\par\smallskip}
% (Parenthesis activation removed: it broke ordinary parentheses in text/math.)

\textbf{B1.}
(a)~$F=Gm_1m_2/r^2=6.67\times10^{-11}\times60\times80/
(0.5)^2=6.67\times10^{-11}\times4800/0.25
=1.28\times10^{-6}\,\text{N}$. \quad
(b)~This is negligible compared with gravitational pull
of Earth on each person ($mg\sim600\,\text{N}$). We do
not feel gravitational attraction to nearby objects
because it is $10^8$ times weaker than Earth's pull. \quad
(c)~If separation doubled: $F'=F/4
=3.2\times10^{-7}\,\text{N}$.

\textbf{B2.}
(a)~$g_{surface}=GM/R^2$;
$g_{orbit}=GM/(R+h)^2=g_{surface}R^2/(R+h)^2
=9.81\times(6.4\times10^6)^2/(6.8\times10^6)^2
=9.81\times0.885=8.68\,\text{m\,s}^{-2}$. \quad
(b)~$v=\sqrt{g_{orbit}(R+h)}=\sqrt{8.68\times6.8\times10^6}
=\sqrt{5.90\times10^7}=7683\,\text{m\,s}^{-1}
\approx7.7\,\text{km\,s}^{-1}$. \quad
(c)~$T=2\pi(R+h)/v=2\pi\times6.8\times10^6/7683
=5559\,\text{s}\approx92.7\,\text{min}$.

\textbf{B3.}
(a)~At geostationary radius $r_{GEO}$:
$GM_E/r_{GEO}^2=\omega^2r_{GEO}$;
$r_{GEO}=(GM_E/\omega^2)^{1/3}$.
$\omega=2\pi/(24\times3600)=7.27\times10^{-5}\,\text{rad\,s}^{-1}$.
$GM_E=gR_E^2=9.81\times(6.4\times10^6)^2
=4.02\times10^{14}$.
$r_{GEO}=(4.02\times10^{14}/(7.27\times10^{-5})^2)^{1/3}
=(7.60\times10^{22})^{1/3}=4.23\times10^7\,\text{m}$
$=42\,300\,\text{km}$ from Earth's centre. \quad
(b)~Height above surface $=42\,300-6400=35\,900\,\text{km}$. \quad
(c)~$v=\omega r_{GEO}=7.27\times10^{-5}\times4.23\times10^7
=3075\,\text{m\,s}^{-1}\approx3.1\,\text{km\,s}^{-1}$.

\textbf{B4.}
(a)~$v_{esc}=\sqrt{2GM/R}=\sqrt{2gR}
=\sqrt{2\times9.81\times6.4\times10^6}
=\sqrt{1.256\times10^8}=11\,200\,\text{m\,s}^{-1}
=11.2\,\text{km\,s}^{-1}$. \quad
(b)~For the Moon ($g_M=1.62\,\text{m\,s}^{-2}$,
$R_M=1.74\times10^6\,\text{m}$):
$v_{esc,M}=\sqrt{2\times1.62\times1.74\times10^6}
=2376\,\text{m\,s}^{-1}=2.4\,\text{km\,s}^{-1}$. \quad
(c)~The Moon's escape velocity is lower than the thermal
speeds of light gases (H$_2$, He) at lunar surface
temperatures, so the Moon cannot retain these gases and
has no atmosphere.

\textbf{B5.}
(a)~$T^2=4\pi^2r^3/(GM_S)$;
$GM_S=4\pi^2r_E^3/T_E^2
=4\times9.87\times(1.5\times10^{11})^3/(3.156\times10^7)^2
=4\times9.87\times3.375\times10^{33}/9.96\times10^{14}
=1.33\times10^{20}$. \quad
$M_S=1.33\times10^{20}/(6.67\times10^{-11})
=2.0\times10^{30}\,\text{kg}$. \quad
(b)~For Mars: $r_M=1.524\,\text{AU}=2.279\times10^{11}\,\text{m}$.
$T_M=T_E(r_M/r_E)^{3/2}=(1.524)^{3/2}=1.881$ years. \quad
(c)~Kepler's Third Law is a consequence of Newton's law
of gravitation and circular orbit dynamics combined.

% ============================================================
%  CHAPTER 9
% ============================================================
\endgroup

\AnswersChapterBand{9}{Elasticity}

\AnswersSectionHeading{Section A}
\begin{multicols}{2}
\begingroup
\let\textbf\AnswersTextbf
\def\quad{ --- }
\setlength{\parindent}{0pt}

\textbf{A1.~B} \quad $E=\tfrac{1}{2}kx^2
=\tfrac{1}{2}\times40\times(0.05)^2=0.05\,\text{J}$.

\textbf{A2.~C} \quad Series: $1/k_s=1/k+1/k=2/k$;
$k_s=k/2=30\,\text{N\,m}^{-1}$.

\textbf{A3.~B} \quad $E=FL/(Ax)
=100\times2/(2\times10^{-6}\times5\times10^{-4})
=200/(10^{-9})=2\times10^{11}\,\text{Pa}$.

\textbf{A4.~C} \quad Brittle: fractures with little or
no plastic deformation (e.g.\ glass, ceramics).

\textbf{A5.~C} \quad Gradient of linear region of
stress-strain graph $=$ stress/strain $=$ Young's
modulus.

\textbf{A6.~B} \quad Strain $=\Delta L/L_0
=(2L-L)/L=1.0$.

\textbf{A7.~C} \quad $R=\rho L/A$; $A\propto d^2$;
double $d\Rightarrow A$ quadruples.
Extension $x=FL/(AE)$; if $A$ quadruples, $x$
becomes $1/4$ of original.

\textbf{A8.~C} \quad Beyond elastic limit: permanent
(plastic) deformation; wire does not return to original
length.

\textbf{A9.~B} \quad $E=\tfrac{1}{2}kx^2\propto x^2$.
Double $x\Rightarrow E$ becomes $4E$.

\textbf{A10.~D} \quad Parallel: $k_p=k_1+k_2+k_3
=90+90+90=270\,\text{N\,m}^{-1}$.

\endgroup
\end{multicols}

\AnswersSectionHeading{Section B}
\begingroup
\let\textbf\AnswersTextbf
\setlength{\parindent}{0pt}
\def\quad{\par\smallskip}
% (Parenthesis activation removed: it broke ordinary parentheses in text/math.)

\textbf{B1.}
(a)~$F=kx=200\times0.06=12\,\text{N}$; this equals the
weight of the mass $\Rightarrow m=12/10=1.2\,\text{kg}$. \quad
(b)~$E=\tfrac{1}{2}kx^2=\tfrac{1}{2}\times200\times0.0036
=0.36\,\text{J}$. \quad
(c)~When released: $0.36=\tfrac{1}{2}mv^2
\Rightarrow v=\sqrt{2\times0.36/1.2}=0.775\,\text{m\,s}^{-1}$.
\quad
(d)~The mass oscillates with $\omega=\sqrt{k/m}
=\sqrt{200/1.2}=12.9\,\text{rad\,s}^{-1}$,
period $T=0.487\,\text{s}$.

\textbf{B2.}
(a)~$E=FL/(Ax)=50\times2.0/(0.785\times10^{-6}
\times3.2\times10^{-3})$. \\
$A=\pi(0.5\times10^{-3})^2=7.85\times10^{-7}\,\text{m}^2$.
$E=100/(2.512\times10^{-9})=3.98\times10^{10}\,\text{Pa}
\approx4\times10^{10}\,\text{Pa}$. \quad
(b)~Stress $=F/A=50/7.85\times10^{-7}
=6.37\times10^7\,\text{Pa}$. \quad
(c)~Strain $=x/L=3.2\times10^{-3}/2.0
=1.6\times10^{-3}$ (dimensionless).

\textbf{B3.}
(a)~Two springs in series: $k_s=k/2=150\,\text{N\,m}^{-1}$.
Extension $=F/k_s=80/150=0.533\,\text{m}$. \quad
(b)~Two in parallel: $k_p=2k=600\,\text{N\,m}^{-1}$.
Extension $=80/600=0.133\,\text{m}$. \quad
(c)~Energy in series: $\tfrac{1}{2}\times150\times0.533^2
=21.3\,\text{J}$. Energy in parallel:
$\tfrac{1}{2}\times600\times0.133^2=5.33\,\text{J}$.
Series stores $4\times$ more energy for the same force.

\textbf{B4.}
(a)~From graph: limit of proportionality at about
$200\,\text{N}$ (where graph starts to curve);
elastic limit slightly beyond this. \quad
(b)~Young's modulus $=$ gradient of linear region
$=\sigma/\varepsilon = (F/A)/(\Delta L/L)$. \quad
(c)~Ductile: large plastic deformation before fracture
(copper, steel); brittle: fractures abruptly (glass,
cast iron). \quad
(d)~Hysteresis: area enclosed by loading and unloading
curves $=$ energy dissipated as heat per cycle.

\textbf{B5.}
(a)~$L_1=2.000\,\text{m}$; $\Delta L_1=F/(EA/L)
=100\times2/(2\times10^{11}\times7.85\times10^{-7})
=200/(1.57\times10^5)=1.27\times10^{-3}\,\text{m}$.
\quad
(b)~$L_2=1.000\,\text{m}$, same $F$, $A$:
$\Delta L_2=100\times1/(E\times A)
=0.637\times10^{-3}\,\text{m}$.
\quad
(c)~Combined: wires in series $\Rightarrow$
total extension $=\Delta L_1+\Delta L_2
=1.91\times10^{-3}\,\text{m}=1.91\,\text{mm}$. \quad
(d)~Effective spring constant:
$k_{eff}=F/\Delta L_{total}=100/0.00191
=52\,400\,\text{N\,m}^{-1}$.

% ============================================================
%  CHAPTER 10
% ============================================================
\endgroup

\AnswersChapterBand{10}{Simple Machines}

\AnswersSectionHeading{Section A}
\begin{multicols}{2}
\begingroup
\let\textbf\AnswersTextbf
\def\quad{ --- }
\setlength{\parindent}{0pt}

\textbf{A1.~B} \quad $MA=\eta\times VR=0.75\times6=4.5$.

\textbf{A2.~C} \quad Third-class lever: effort applied
between fulcrum and load; effort arm $<$ load arm;
$MA<1$ always.

\textbf{A3.~C} \quad $VR=5$; ideal $MA=VR=5$.
Effort $=L/MA=1000/5=200\,\text{N}$.

\textbf{A4.~B} \quad $VR=2\pi L/p
=2\pi\times0.20/0.004=100\pi=314$.

\textbf{A5.~B} \quad $VR=\text{length}/\text{height}
=8/2=4$.

\textbf{A6.~B} \quad Ideal wheel-and-axle:
Effort$\times$wheel radius $=$ Load$\times$axle radius.
$E=600\times8/40=120\,\text{N}$.

\textbf{A7.~B} \quad An ideal machine has efficiency
$=100\%$ ($MA=VR$). No machine creates energy; $MA$
can never exceed $VR$ in a real machine.

\textbf{A8.~C} \quad $MA=$ effort arm/load arm.
Wheelbarrow: load is $0.3\,\text{m}$ from wheel
(fulcrum); handles are $1.5\,\text{m}$ from load,
so effort arm $=0.3+1.5=1.8\,\text{m}$ from fulcrum.
Actually for a second-class lever:
$MA=$ effort arm$/$load arm $=1.8/0.3=6$... hmm.

Let me re-read: ``wheel (fulcrum) $0.3\,\text{m}$ from the
load and handles $1.5\,\text{m}$ from the load.''
Load arm $=0.3\,\text{m}$.
Effort arm $=$ distance from fulcrum to effort
$=0.3+1.5=1.8\,\text{m}$.
$MA=1.8/0.3=6$. But option C $=5.0$ and D $=6.0$.
$MA=6.0$ $\Rightarrow$ \textbf{D}.

(Note to Falcon: on proofing verify whether the
wheelbarrow problem intends effort arm $=1.5\,\text{m}$
from the fulcrum or $1.5\,\text{m}$ from the load.
If effort arm $=1.5\,\text{m}$ from fulcrum:
$MA=1.5/0.3=5.0$ $\Rightarrow$ C.)

\textbf{A9.~C} \quad Speed ratio $=$ driver teeth/driven
teeth $=20/80=1/4$.
Driven speed $=400\times(1/4)=100\,\text{rpm}$.

\textbf{A10.~C} \quad $L=\eta\times VR\times E
=0.60\times5\times200=600\,\text{N}$.

\endgroup
\end{multicols}

\AnswersSectionHeading{Section B}
\begingroup
\let\textbf\AnswersTextbf
\setlength{\parindent}{0pt}
\def\quad{\par\smallskip}
% (Parenthesis activation removed: it broke ordinary parentheses in text/math.)

\textbf{B1.}
(a)~$VR=\text{effort distance}/\text{load distance}
=1.2/0.4=3$; $\eta=MA/VR$; $MA=L/E=300/120=2.5$;
$\eta=2.5/3=83.3\%$. \quad
(b)~Energy input $=$ Effort$\times$effort distance
$=120\times1.2=144\,\text{J}$.
Useful output $=300\times0.4=120\,\text{J}$.
Energy wasted $=24\,\text{J}$ (friction in pivots). \quad
(c)~Frictional force (converted): work lost/effort
distance $=24/1.2=20\,\text{N}$.

\textbf{B2.}
(a)~$VR=4$ (four load-bearing rope sections).
Ideal effort $=800/4=200\,\text{N}$. \quad
(b)~Actual effort with $\eta=0.80$:
$E=L/(\eta\times VR)=800/(0.80\times4)=250\,\text{N}$.
\quad
(c)~Work input $=E\times d_E$; work output
$=L\times d_L$; $d_E=VR\times d_L=4\times2=8\,\text{m}$.
Work input $=250\times8=2000\,\text{J}$.
Work output $=800\times2=1600\,\text{J}$.
$\eta=1600/2000=80\%$. \checkmark

\textbf{B3.}
(a)~For $n$ moveable pulleys: $VR=2n$.
$n=3\Rightarrow VR=6$. \quad
(b)~Ideal effort $=1500/6=250\,\text{N}$. \quad
(c)~With $\eta=0.75$: actual effort
$=1500/(0.75\times6)=333\,\text{N}$. \quad
(d)~To lift $2\,\text{m}$: effort moves $6\times2
=12\,\text{m}$.

\textbf{B4.}
(a)~$VR=$ wheel circumference/axle circumference
$=2\pi\times0.35/(2\pi\times0.05)=7.0$. \quad
(b)~$\eta=MA/VR$; $MA=L/E=490/(0.35\times10^{-1}
\times980)\approx...$
Actually: $E=490/(VR\times\eta)=490/(7\times0.85)
=82.4\,\text{N}$. \quad
(c)~Screwdriver: VR $=2\pi r_{handle}/r_{tip}$;
longer handle, wider grip $\Rightarrow$ higher $VR$,
more torque for same effort.

\textbf{B5.}
(a)~Gears A ($N_A=12$), B ($N_B=48$), C
($N_C=16$): A drives B, B drives C.
Speed ratio A:B $=48/12=4$ (B runs at $n/4$).
B and C on same axle? If so C runs at same speed as B.
If B drives C: speed ratio B:C $=16/48$... without
more detail, general answer: overall ratio
$=N_A\times N_C/(N_B\times...) $ depends on arrangement. \quad
(b)~Torque increases when speed decreases (power $=\tau\omega=$
constant for ideal gear). \quad
(c)~Industrial gearboxes step down speed from motor to
increase torque for rotating heavy loads.

% ============================================================
%  CHAPTER 11
% ============================================================
\endgroup

\AnswersChapterBand{11}{Density, Pressure and Upthrust}

\AnswersSectionHeading{Section A}
\begin{multicols}{2}
\begingroup
\let\textbf\AnswersTextbf
\def\quad{ --- }
\setlength{\parindent}{0pt}

\textbf{A1.~B} \quad $P=F/A=60/0.3=200\,\text{Pa}$.

\textbf{A2.~B} \quad $h=P/(\rho g)
=5\times10^5/(1000\times10)=50\,\text{m}$.

\textbf{A3.~C} \quad Upthrust $=12-8=4\,\text{N}$.

\textbf{A4.~B} \quad Floating: $\rho_{wood}V g
=\rho_{water}(0.6V)g\Rightarrow\rho_{wood}
=600\,\text{kg\,m}^{-3}$.

\textbf{A5.~B} \quad Pascal's principle: pressure
applied to enclosed fluid transmits equally in all
directions $\Rightarrow$ hydraulic systems.

\textbf{A6.~C} \quad Relative density $8.5\times1000
=8500\,\text{kg\,m}^{-3}$.

\textbf{A7.~B} \quad Neutral buoyancy: upthrust $=$
weight $\Rightarrow$ density of object $=$ density of
fluid.

\textbf{A8.~C} \quad $P=\rho gh$; depends only on
depth and fluid density, not container shape or total
mass.

\textbf{A9.~C} \quad $F_2=F_1\times A_2/A_1
=10\times200/2=1000\,\text{N}$.

\textbf{A10.~C} \quad Upthrust (Archimedes' principle)
$=$ weight of fluid displaced; reduces apparent weight
of anchor.

\endgroup
\end{multicols}

\AnswersSectionHeading{Section B}
\begingroup
\let\textbf\AnswersTextbf
\setlength{\parindent}{0pt}
\def\quad{\par\smallskip}
% (Parenthesis activation removed: it broke ordinary parentheses in text/math.)

\textbf{B1.}
(a)~$\rho=m/V=7.9\,\text{g/cm}^3=7900\,\text{kg\,m}^{-3}$.
\quad
(b)~Apparent weight $=W-U=W-\rho_{water}Vg$.
$V=m/\rho_{Fe}=0.500/7900=6.33\times10^{-5}\,\text{m}^3$.
$U=1000\times6.33\times10^{-5}\times10=0.633\,\text{N}$.
Apparent weight $=0.500\times10-0.633=4.37\,\text{N}$. \quad
(c)~In oil ($\rho=800\,\text{kg\,m}^{-3}$):
$U=800\times6.33\times10^{-5}\times10=0.506\,\text{N}$.
Apparent weight $=5-0.506=4.49\,\text{N}$.

\textbf{B2.}
(a)~Weight of ship $=W$; volume of water displaced
$=W/(\rho g)$. If $W=100\times10^6\,\text{N}$:
$V=100\times10^6/(1025\times10)
=9756\,\text{m}^3$. \quad
(b)~Extra $2000\,\text{kg}$ cargo increases total weight;
ship sinks deeper until extra upthrust $=2000\times10
=20\,000\,\text{N}$ more. Extra volume
$=20\,000/(1025\times10)=1.95\,\text{m}^3$.
If cross-section $\approx500\,\text{m}^2$:
extra depth $\approx0.004\,\text{m}=4\,\text{mm}$.

\textbf{B3.}
(a)~Total pressure at depth $h=60\,\text{m}$:
$P=P_{atm}+\rho gh=101\,000+1025\times10\times60
=101\,000+615\,000=716\,000\,\text{Pa}$. \quad
(b)~Force on porthole ($d=0.25\,\text{m}$, $A=\pi r^2
=0.0491\,\text{m}^2$):
$F=P\times A=716\,000\times0.0491=35\,150\,\text{N}
\approx35\,\text{kN}$. \quad
(c)~At SCUBA depth the diver breathes compressed air at
ambient pressure; the high nitrogen partial pressure
causes narcosis at deep depths.

\textbf{B4.}
(a)~Small piston area $A_1=\pi(0.02)^2
=1.257\times10^{-3}\,\text{m}^2$.
Large piston $A_2=\pi(0.15)^2=7.07\times10^{-2}\,\text{m}^2$.
$F_2=F_1\times A_2/A_1=500\times7.07/1.257\times10^{-1}
=500\times56.25=28\,125\,\text{N}$. \quad
(b)~$VR=A_2/A_1=56.25$; to lift $0.5\,\text{m}$,
small piston must move $56.25\times0.5=28.1\,\text{m}$.
\quad
(c)~$\eta=W_{out}/W_{in}=28125\times0.5/(500\times28.1)
=14062.5/14050=100\%$ (ideal). In practice
$<100\%$ due to friction and leakage.

\textbf{B5.}
(a)~$RD=\rho_{substance}/\rho_{water}$. Measure weight
$W_1$ in air, weight $W_2$ when fully submerged in
water. $RD=W_1/(W_1-W_2)$. \quad
(b)~For the solid: $W_1=2.5\,\text{N}$;
$W_2=1.8\,\text{N}$; $RD=2.5/0.7=3.57$. \quad
(c)~$\rho=3.57\times1000=3570\,\text{kg\,m}^{-3}$
(consistent with aluminium alloy or similar material). \quad
(d)~For a floating body: $\rho_{object}\times V_{total}
=\rho_{fluid}\times V_{submerged}$;
$RD=V_{submerged}/V_{total}=$ fraction submerged.

% ============================================================
%  CHAPTER 12
% ============================================================
\endgroup

\AnswersChapterBand{12}{Surface Tension and Viscosity}

\AnswersSectionHeading{Section A}
\begin{multicols}{2}
\begingroup
\let\textbf\AnswersTextbf
\def\quad{ --- }
\setlength{\parindent}{0pt}

\textbf{A1.~B} \quad Surface molecules have a net
inward cohesive force (no molecules above them to
balance the downward attraction).

\textbf{A2.~C} \quad Soap bubble has two liquid
surfaces: excess pressure $=4\gamma/r$.

\textbf{A3.~C} \quad $h=2\gamma\cos\theta/(\rho gr)$.
If $r\to r/3$: $h\to3h$.

\textbf{A4.~B} \quad Mercury's cohesion exceeds its
adhesion to glass; it forms a convex meniscus and
does not wet glass.

\textbf{A5.~C} \quad Surfactants (soap) reduce surface
tension by disrupting hydrogen bonding at the water
surface.

\textbf{A6.~C} \quad Stokes' law $F=6\pi\eta rv$
applies only for laminar (streamline) flow around
the sphere.

\textbf{A7.~C} \quad At terminal velocity net force
$=0$: weight $=$ upthrust $+$ drag
($mg=U+6\pi\eta rv$).

\textbf{A8.~B} \quad From Stokes' terminal velocity:
$v_t=2r^2(\rho_s-\rho_f)g/(9\eta)\propto r^2$.

\textbf{A9.~B} \quad Viscosity of liquids decreases
with temperature (molecules have more thermal energy
to overcome intermolecular binding).

\textbf{A10.~B} \quad Pressure inside a bubble
$\propto1/r$; smaller bubble has higher pressure;
air flows from smaller to larger bubble.

\endgroup
\end{multicols}

\AnswersSectionHeading{Section B}
\begingroup
\let\textbf\AnswersTextbf
\setlength{\parindent}{0pt}
\def\quad{\par\smallskip}
% (Parenthesis activation removed: it broke ordinary parentheses in text/math.)

\textbf{B1.}
(a)~$\gamma=F/(2L)=0.015/(2\times0.1)
=0.075\,\text{N\,m}^{-1}$ (factor of 2 because frame
film has two surfaces). \quad
(b)~$h=2\gamma\cos\theta/(\rho gr)
=2\times0.073\times1/(1000\times10\times0.5\times10^{-3})
=0.146/5=0.0292\,\text{m}=2.92\,\text{cm}$. \quad
(c)~Adding detergent reduces $\gamma$; $h$ decreases
proportionally.

\textbf{B2.}
(a)~Weight $=mg=\tfrac{4}{3}\pi r^3(\rho_s-\rho_f)g$.
Stokes drag $=6\pi\eta r v_t$.
At terminal velocity:
$v_t=\frac{2r^2(\rho_s-\rho_f)g}{9\eta}$. \quad
(b)~$r=0.5\times10^{-3}\,\text{m}$,
$\rho_s=7800$, $\rho_f=900$, $\eta=0.5\,\text{Pa\,s}$,
$g=10$:
$v_t=2\times(0.5\times10^{-3})^2\times6900\times10/
(9\times0.5)=2\times2.5\times10^{-7}\times69000/4.5
=3.45\times10^{-2}/4.5=7.67\times10^{-3}\,\text{m\,s}^{-1}$. \quad
(c)~Reynolds number $Re=\rho_f v_t(2r)/\eta
=900\times7.67\times10^{-3}\times10^{-3}/0.5=0.0138$.
$Re\ll1\Rightarrow$ laminar flow. Stokes' law valid.

\textbf{B3.}
(a)~Excess pressure $P=2\gamma/r
=2\times0.025/0.01=5\,\text{Pa}$ (water droplet, one
surface). \quad
(b)~Soap bubble: $P=4\gamma/r=4\times0.025/0.005
=20\,\text{Pa}$. \quad
(c)~When two bubbles of different radii merge: the
smaller (higher pressure) pushes into the larger;
the final bubble has radius determined by volume
conservation and pressure equalisation. Intermediate
surface becomes convex toward the larger bubble.

\textbf{B4.}
(a)~$\eta=F/(Av/d)$; units $=\text{N}/(\text{m}^2
\times\text{m\,s}^{-1}/\text{m})=\text{N\,s\,m}^{-2}
=\text{Pa\,s}$. \quad
(b)~$F=\eta A v/d=1.0\times10^{-3}\times0.01
\times1.0/5\times10^{-4}=2\times10^{-2}\,\text{N}
=0.02\,\text{N}$. \quad
(c)~As temperature rises, water viscosity falls
(from $\sim1\,\text{mPa\,s}$ at $20°\text{C}$ to
$\sim0.3\,\text{mPa\,s}$ at $80°\text{C}$).
The engine runs more smoothly because oil also becomes
less viscous when hot, but too low viscosity causes
inadequate lubrication --- hence motor oil grades
(5W-30 etc.) are designed to maintain viscosity across
temperature ranges.

\textbf{B5.}
(a)~$h=2\gamma/(\rho g r)$ where $r$ is tube radius.
Measure $h$ for several known radii $r$; plot $h$ vs
$1/r$; gradient $=2\gamma/(\rho g)$; hence $\gamma$. \quad
(b)~Sources of error: contamination of tube, measuring
height with ruler, meniscus shape uncertainty,
temperature variation. \quad
(c)~Insects walking on water: their legs dimple the
surface film; weight $\leq$ surface tension force;
$mg\leq2\gamma L$ where $L$ is total contact
perimeter.

% ============================================================
%  CHAPTER 13
% ============================================================
\endgroup

\AnswersChapterBand{13}{Temperature and Thermal Expansion}

\AnswersSectionHeading{Section A}
\begin{multicols}{2}
\begingroup
\let\textbf\AnswersTextbf
\def\quad{ --- }
\setlength{\parindent}{0pt}

\textbf{A1.~A} \quad $T(\text{K})=-40+273=233\,\text{K}$.

\textbf{A2.~B} \quad Linear interpolation:
$L_{25}=L_0+(L_{100}-L_0)\times25/100
=2+(10\times0.25)=2+2.5=4.5\,\text{cm}$.

\textbf{A3.~A} \quad $\Delta L=\alpha L_0\Delta T
=12\times10^{-6}\times2.0\times50=1.2\times10^{-3}\,\text{m}$.
New length $=2.0+0.0012=2.0012\,\text{m}$.

\textbf{A4.~C} \quad $\gamma=3\alpha$.

\textbf{A5.~B} \quad Water is densest at $4°\text{C}$
due to the hydrogen-bond structure of ice forming at
lower temperatures.

\textbf{A6.~B} \quad Different $\alpha$ values cause
one strip to expand more than the other; the strip
curves toward the less-expanding side.

\textbf{A7.~B} \quad Apparent expansion
$=$ real expansion $-$ expansion of container.

\textbf{A8.~C} \quad Kelvin scale: $0\,\text{K}
=$ absolute zero.

\textbf{A9.~B} \quad The hole expands exactly as if
the removed material were still there; it gets larger.

\textbf{A10.~B} \quad Zeroth Law: if A and B are each
in thermal equilibrium with C, then A and B are in
thermal equilibrium with each other. This defines
temperature and justifies thermometers.

\endgroup
\end{multicols}

\AnswersSectionHeading{Section B}
\begingroup
\let\textbf\AnswersTextbf
\setlength{\parindent}{0pt}
\def\quad{\par\smallskip}
% (Parenthesis activation removed: it broke ordinary parentheses in text/math.)

\textbf{B1.}
(a)~$\Delta L=\alpha L_0\Delta T
=11\times10^{-6}\times50\times80=0.044\,\text{m}$. \quad
(b)~$\Delta V=\gamma V_0\Delta T=3\alpha V_0\Delta T
=3\times11\times10^{-6}\times0.01\times80
=2.64\times10^{-5}\,\text{m}^3$. \quad
(c)~Gap needed $=\Delta L=0.044\,\text{m}=44\,\text{mm}$
per $50\,\text{m}$. Provide expansion joints every
$10\,\text{m}$ or at sufficient intervals.

\textbf{B2.}
(a)~$\Delta L=\alpha L\Delta T=12\times10^{-6}
\times1.0\times200=2.4\times10^{-3}\,\text{m}=2.4\,\text{mm}$. \quad
(b)~If both ends are fixed, the rod cannot expand and
develops a compressive stress.
$\sigma=E\alpha\Delta T=2\times10^{11}
\times12\times10^{-6}\times200=4.8\times10^8\,\text{Pa}$. \quad
(c)~This stress may exceed the yield strength of the
material, causing buckling. Rail tracks buckle in
extreme heat if expansion gaps are insufficient.
(This is a genuine hazard on Nigerian railways during
the dry season heat.)

\textbf{B3.}
(a)~$\Delta V_{real}=\gamma_{Hg}V_0\Delta T
=1.82\times10^{-4}\times V_0\times\Delta T$.
$\Delta V_{apparent}=(\gamma_{Hg}-\gamma_{glass})
V_0\Delta T
=(1.82-0.027)\times10^{-4}V_0\Delta T$. \quad
(b)~Apparent coefficient $\gamma_{app}
=\gamma_{real}-\gamma_{container}$. \quad
(c)~Correction factor $=\gamma_{glass}\times V_0\times\Delta T$. \quad
(d)~Mercury thermometer accurate to $\pm0.1°\text{C}$ with
careful calibration; alcohol thermometer suitable for
lower temperatures (alcohol freezes at $-115°\text{C}$
vs mercury at $-39°\text{C}$).

\textbf{B4.}
(a)~$\Delta L/L=\alpha\Delta T$;
$\alpha=\Delta L/(L\Delta T)
=1.0\times10^{-3}/(2.0\times100)
=5.0\times10^{-6}\,\text{K}^{-1}$.
Likely material: invar or similar low-expansion alloy. \quad
(b)~$\Delta V=3\alpha V_0\Delta T
=3\times5\times10^{-6}\times1.5\times100
=2.25\times10^{-3}\,\text{m}^3$. \quad
(c)~At $120°\text{C}$: $\Delta L=5\times10^{-6}
\times2.0\times100=10^{-3}\,\text{m}$.

\textbf{B5.}
(a)~Thermostat: bimetallic strip bends on heating;
when strip touches contact, circuit opens (or closes)
switching off (or on) the heater. \quad
(b)~Strip A expands more (higher $\alpha$); strip curves
toward B side. \quad
(c)~Temperature range: set by the distance between the
strip and the contact; adjusting a screw changes this
distance. \quad
(d)~Applications: electric iron, oven, fridge, fire
alarm, kettle.

% ============================================================
%  CHAPTER 14
% ============================================================
\endgroup

\AnswersChapterBand{14}{Gas Laws and Kinetic Theory}

\AnswersSectionHeading{Section A}
\begin{multicols}{2}
\begingroup
\let\textbf\AnswersTextbf
\def\quad{ --- }
\setlength{\parindent}{0pt}

\textbf{A1.~C} \quad $P_1V_1=P_2V_2$:
$100\times8=P_2\times2\Rightarrow P_2=400\,\text{kPa}$.

\textbf{A2.~C} \quad $V\propto T$ (Charles' Law).
$T_1=300\,\text{K}$; double $V\Rightarrow T_2=600\,\text{K}
=327°\text{C}$.

\textbf{A3.~B} \quad Boyle's Law: $P\propto1/V$, so $P$
vs $1/V$ is a straight line through the origin.

\textbf{A4.~C} \quad $\overline{KE}=\tfrac{3}{2}k_BT$;
depends only on absolute temperature.

\textbf{A5.~C} \quad $c_{rms}\propto\sqrt{T}$.
Double $c_{rms}\Rightarrow T$ quadruples.

\textbf{A6.~B} \quad $P=nRT/V
=2\times8.314\times300/(50\times10^{-3})
=4988/0.05=99\,760\approx99.8\,\text{kPa}$.

\textbf{A7.~C} \quad Dalton's Law: $P_{total}
=\sum P_i$.

\textbf{A8.~B} \quad Ideal gas: zero molecular volume
and no intermolecular forces.

\textbf{A9.~B} \quad Molar volume at STP $\approx22.4\,\text{L}$.

\textbf{A10.~C} \quad At high pressure, molecules are
packed closely; their own volume becomes a significant
fraction of total volume.

\endgroup
\end{multicols}

\AnswersSectionHeading{Section B}
\begingroup
\let\textbf\AnswersTextbf
\setlength{\parindent}{0pt}
\def\quad{\par\smallskip}
% (Parenthesis activation removed: it broke ordinary parentheses in text/math.)

\textbf{B1.}
(a)~At surface: $V_s=nRT/P=n\times8.314\times293/
(1.013\times10^5)$. Volume of tank: $12\,\text{L}$ at
$200\,\text{atm}$; at surface pressure:
$V=12\times200=2400\,\text{L}=2.4\,\text{m}^3$.
\quad
(b)~At $30\,\text{m}$ ($P=4.07\,\text{atm}$,
$T=283\,\text{K}$), using $P_1V_1/T_1=P_2V_2/T_2$:
$V_2=P_1V_1T_2/(T_1P_2)=200\times12\times283/
(293\times4.07)=678\,600/1192=569\,\text{L}$. \quad
(c)~Each breath uses the same volume; at $30\,\text{m}$
the air in that volume is at 4.07 atm, containing
$4.07\times$ as many molecules. Diver uses $4.07\times$
as much air per breath $\Rightarrow$ tank lasts
$2400/(4.07\times V_{breath})$ breaths vs
$2400/V_{breath}$ at surface.
Ratio $=1/4.07$; diver breathes $4\times$ shorter.

\textbf{B2.}
(a)~$n=PV/(RT)=8.0\times10^5\times0.015/
(8.314\times298)=12000/2477.6=4.84\,\text{mol}$. \quad
(b)~$m=n\times M=4.84\times44=213\,\text{g}$. \quad
(c)~At $50°\text{C}=323\,\text{K}$:
$P'=P\times T'/T=8.0\times10^5\times323/298
=8.67\times10^5\,\text{Pa}<9.5\times10^5\,\text{Pa}$.
Valve does not open. \quad
(d)~At $300°\text{C}=573\,\text{K}$:
$P'=8.0\times10^5\times573/298=1.54\times10^6\,\text{Pa}$.
Far above relief valve pressure. Dangerous: cylinder may
explode.

\textbf{B3.}
(a)~$PV=\tfrac{1}{3}Nm\langle c^2\rangle$ and
$PV=Nk_BT$; equating:
$\tfrac{1}{3}Nm\langle c^2\rangle=Nk_BT\Rightarrow
\tfrac{1}{2}m\langle c^2\rangle=\tfrac{3}{2}k_BT$.
\checkmark \quad
(b)~At $300\,\text{K}$:
$\overline{KE}=\tfrac{3}{2}\times1.38\times10^{-23}
\times300=6.21\times10^{-21}\,\text{J}$.
At $600\,\text{K}$: $1.24\times10^{-20}\,\text{J}$. \quad
(c)~$c_{rms}=\sqrt{3RT/M}$.
At $300\,\text{K}$:
$c=\sqrt{3\times8.314\times300/0.028}=517\,\text{m\,s}^{-1}$.
At $600\,\text{K}$:
$c=\sqrt{3\times8.314\times600/0.028}=730\,\text{m\,s}^{-1}$.
\quad
(d)~$U=\tfrac{3}{2}nRT=\tfrac{3}{2}\times2\times8.314
\times300=7483\,\text{J}\approx7.48\,\text{kJ}$.

\textbf{B4.}
(a)~$P_i=n_iRT/V$:
$P_{He}=1.0\times8.314\times300/0.020=124\,710\,\text{Pa}
\approx125\,\text{kPa}$;
$P_{Ne}=0.5\times8.314\times300/0.020=62\,355\,\text{Pa}
\approx62.4\,\text{kPa}$;
$P_{Ar}=1.5\times8.314\times300/0.020=187\,065\,\text{Pa}
\approx187\,\text{kPa}$.
\quad
(b)~$P_{total}=125+62.4+187=374\,\text{kPa}$. \quad
(c)~$n_{total}=3.0\,\text{mol}$;
$P=nRT/V=3\times8.314\times300/0.020=374\,130\,\text{Pa}
\approx374\,\text{kPa}$. \checkmark \quad
(d)~Without Ar ($n=1.5\,\text{mol}$):
$P=1.5\times8.314\times300/0.020=187\,\text{kPa}$. \quad
(e)~$c_{rms}\propto1/\sqrt{M}$.
Ratio $=\sqrt{M_{Ar}/M_{He}}=\sqrt{40/4}=\sqrt{10}
=3.16$.
Helium molecules move $3.16\times$ faster than argon
at the same temperature.

\textbf{B5.}
(a)~$PV=nRT\Rightarrow P(M/\rho)=nRT$. For $n=1$
mole, mass $=M$, volume $=M/\rho$:
$P(M/\rho)=RT\Rightarrow\rho=PM/(RT)$. \checkmark \quad
(b)~$\rho=1.013\times10^5\times0.029/(8.314\times293)
=2938/2436=1.206\,\text{kg\,m}^{-3}$. \quad
(c)~For $\rho'=0.80\rho$: $T'=T/0.80=293/0.80
=366\,\text{K}=93°\text{C}$. \quad
(d)~At $35°\text{C}=308\,\text{K}$:
$\rho=1.01\times10^5\times0.029/(8.314\times308)
=1.148\,\text{kg\,m}^{-3}$.
Less dense than at $20°\text{C}$; less air entering
engines per cycle; less lift generated per unit area of
wing at same speed $\Rightarrow$ longer take-off run
needed.

% ============================================================
%  CHAPTER 15
% ============================================================
\endgroup

\AnswersChapterBand{15}{Heat Energy and Calorimetry}

\AnswersSectionHeading{Section A}
\begin{multicols}{2}
\begingroup
\let\textbf\AnswersTextbf
\def\quad{ --- }
\setlength{\parindent}{0pt}

\textbf{A1.~C} \quad $Q=mc\Delta\theta
=2\times390\times50=39\,000\,\text{J}$.

\textbf{A2.~C} \quad During melting of a pure substance,
temperature remains constant (latent heat absorbed
breaks bonds without raising temperature).

\textbf{A3.~B} \quad Vaporisation must completely
overcome all intermolecular forces and do work against
atmospheric pressure; melting only partially disrupts
the structure.

\textbf{A4.~C} \quad $Q=Pt=500\times240=120\,000\,\text{J}$.
$c=Q/(m\Delta\theta)=120\,000/(2\times30)
=2000\,\text{J\,kg}^{-1}\text{K}^{-1}$.

\textbf{A5.~A} \quad Heat available from hot water to
cool to $0°\text{C}$:
$Q=0.1\times4200\times80=33\,600\,\text{J}$.
Heat to melt all ice: $0.1\times3.36\times10^5
=33\,600\,\text{J}$.
Exactly equal $\Rightarrow$ all ice melts and final
temperature $=0°\text{C}$.

\textbf{A6.~B} \quad Evaporation occurs at the surface
at any temperature and removes the fastest molecules,
cooling the liquid.

\textbf{A7.~B} \quad Heat lost to surroundings means
less heat absorbed by the substance; temperature rise
is less than expected; $c=Q_{in}/(m\Delta\theta)$ gives
a value larger than true $c$.

\textbf{A8.~C} \quad $Q=mL_v=0.200\times2.26\times10^6
=452\,000\,\text{J}$.

\textbf{A9.~B} \quad Method of mixtures uses
conservation of energy: heat lost by hot object
$=$ heat gained by cool object.

\textbf{A10.~C} \quad Steam at $100°\text{C}$ releases
$2.26\times10^6\,\text{J\,kg}^{-1}$ of latent heat
when it condenses on skin, causing much more severe
burns than liquid water at the same temperature.

\endgroup
\end{multicols}

\AnswersSectionHeading{Section B}
\begingroup
\let\textbf\AnswersTextbf
\setlength{\parindent}{0pt}
\def\quad{\par\smallskip}
% (Parenthesis activation removed: it broke ordinary parentheses in text/math.)

\textbf{B1.}
(a)~$Q=mc\Delta\theta=5.0\times4200\times85
=1\,785\,000\,\text{J}$.
Time $=Q/P=1\,785\,000/2000=892.5\,\text{s}
\approx14.9\,\text{min}$. \quad
(b)~$Q_{vap}=1.0\times2.26\times10^6=2\,260\,000\,\text{J}$.
Time $=2\,260\,000/2000=1130\,\text{s}\approx18.8\,\text{min}$.
\quad
(c)~If actual time is $12\%$ longer:
$\eta=1/1.12\approx89\%$. \quad
(d)~Losses include: heat to container walls; heat
radiated/convected to surroundings; heat carried away
by steam.

\textbf{B2.}
(a)~Heat gained by water:
$Q_w=0.300\times4200\times10=12\,600\,\text{J}$.
Heat gained by calorimeter:
$Q_c=0.150\times900\times10=1350\,\text{J}$.
Total heat gained $=13\,950\,\text{J}$. \quad
(b)~Heat lost by steel $=13\,950\,\text{J}$ (assuming
perfect insulation). \quad
(c)~$c_{steel}=Q/(m\Delta\theta)
=13\,950/(0.200\times92)=13\,950/18.4
=758\,\text{J\,kg}^{-1}\text{K}^{-1}$. \quad
(d)~This is higher than the accepted value of
$450\,\text{J\,kg}^{-1}\text{K}^{-1}$ because heat
was lost to surroundings; the measured $\Delta\theta$
is smaller than ideal, making calculated $c$ larger.

\textbf{B3.}
(a)~Heat to bring ice to $0°\text{C}$:
$Q_1=0.400\times2100\times15=12\,600\,\text{J}$. \quad
(b)~Heat from warm water cooling to $0°\text{C}$:
$Q_{avail}=0.600\times4200\times50=126\,000\,\text{J}$.
\quad
(c)~Heat to melt all ice:
$Q_{melt}=0.400\times3.36\times10^5=134\,400\,\text{J}$.
Heat remaining after warming ice to $0°\text{C}$:
$126\,000-12\,600=113\,400\,\text{J}<134\,400\,\text{J}$.
\textbf{Not all ice melts.} Some ice remains. \quad
(d)~Since not all ice melts: temperature stays at
$0°\text{C}$.
Mass of ice melted:
$m_{melted}=113\,400/(3.36\times10^5)=0.337\,\text{kg}$.
Unmelted ice: $0.400-0.337=0.063\,\text{kg}=63\,\text{g}$.
Final temperature $=0°\text{C}$.

\textbf{B4.}
(a)~$Q_1=mc\Delta\theta=0.500\times4200\times25
=52\,500\,\text{J}$ (cooling water from $25°\text{C}$
to $0°\text{C}$).
$Q_2=mL_f=0.500\times3.36\times10^5=168\,000\,\text{J}$
(freezing at $0°\text{C}$).
$Q_3=mc_{ice}\Delta\theta=0.500\times2100\times10
=10\,500\,\text{J}$ (cooling ice from $0°\text{C}$
to $-10°\text{C}$).
Total $=231\,000\,\text{J}$. \quad
(b)~Time $=231\,000/80=2888\,\text{s}\approx48\,\text{min}$.
\quad
(c)~Sketch: temperature drops from $25°\text{C}$ to
$0°\text{C}$ (sloped), then flat plateau at $0°\text{C}$
(freezing, longest section), then drops to $-10°\text{C}$.
\quad
(d)~The plateau (latent heat phase) requires
$168\,000\,\text{J}$ versus only $63\,000\,\text{J}$
for all the sensible heat; the latent heat plateau
is far longer than either cooling section.

\textbf{B5.}
(a)~(i)~Fastest molecules escape the surface; average
speed of remaining molecules falls; lower average KE
means lower temperature.
(ii)~At high altitude, atmospheric pressure is lower;
the boiling point is the temperature at which vapour
pressure equals atmospheric pressure; lower pressure
$\Rightarrow$ lower boiling point.
(iii)~Pressure cooker: higher pressure $\Rightarrow$
higher temperature needed to reach boiling
(vapour pressure must equal higher external pressure).
\quad
(b)~$Q=mc\Delta\theta+mL_v
=0.300\times4200\times91+0.300\times2.26\times10^6
=114\,660+678\,000=792\,660\,\text{J}\approx793\,\text{kJ}$.
\quad
(c)~$Q_{removed}=1.5\times2.4\times10^6
=3.6\times10^6\,\text{J}$.
If not sweating: $\Delta\theta=Q/(mc_{body})
=3.6\times10^6/(70\times3500)=14.7°\text{C}$.
Body temperature would rise from $37°\text{C}$
to $51.7°\text{C}$ --- fatal. Sweating is
essential for thermoregulation.

% ============================================================
%  CHAPTER 16
% ============================================================
\endgroup

\AnswersChapterBand{16}{Heat Transfer}

\AnswersSectionHeading{Section A}
\begin{multicols}{2}
\begingroup
\let\textbf\AnswersTextbf
\def\quad{ --- }
\setlength{\parindent}{0pt}

\textbf{A1.~C} \quad Only radiation can travel through
vacuum (as electromagnetic waves). Conduction and
convection require a medium.

\textbf{A2.~A} \quad $\dot{Q}=kA\Delta T/L
=50\times(2\times10^{-4})\times80/0.4=2\,\text{W}$.

\textbf{A3.~D} \quad Stefan-Boltzmann: $P=\varepsilon\sigma AT^4$;
power $\propto T^4$.

\textbf{A4.~B} \quad Metals are good thermal conductors;
free electrons rapidly carry heat away from your hand.

\textbf{A5.~B} \quad Land heats faster (lower specific
heat capacity); hot air rises over land; cooler sea air
flows inland as sea breeze.

\textbf{A6.~D} \quad $P\propto T^4$;
$(2T)^4=16T^4\Rightarrow P$ multiplied by $16$.

\textbf{A7.~B} \quad Vacuum prevents conduction (no
material) and convection (no fluid).

\textbf{A8.~B} \quad $R_A=L/(kA)$; $R_B=L/(2kA)$;
$R_{total}=R_A+R_B=L/(kA)+L/(2kA)=3L/(2kA)$.
$R_A/R_{total}=(L/kA)/(3L/2kA)=2/3$.
Ratio $R_A:R_{total}=2:3$.

\textbf{A9.~B} \quad Greenhouse effect: atmosphere
transparent to short-wave solar radiation; absorbs
outgoing long-wave infrared; re-radiates back to Earth.

\textbf{A10.~A} \quad Ratio $=\varepsilon_s/\varepsilon_{bb}
=0.02/1.0=1:50$.
\endgroup
\end{multicols}

\AnswersSectionHeading{Section B}
\begingroup
\let\textbf\AnswersTextbf
\setlength{\parindent}{0pt}
\def\quad{\par\smallskip}
% (Parenthesis activation removed: it broke ordinary parentheses in text/math.)


\textbf{B1.}
(a)~$\dot{Q}=kA\Delta T/L=0.70\times15\times14/0.22
=147/0.22=668\,\text{W}$. \quad
(b)~$R_{brick}=0.22/(0.70\times15)=0.02095\,\text{K\,W}^{-1}$.
$R_{insul}=0.03/(0.038\times15)=0.05263\,\text{K\,W}^{-1}$.
$R_{total}=0.0736\,\text{K\,W}^{-1}$.
$\dot{Q}'=14/0.0736=190\,\text{W}$. \quad
(c)~Reduction $=(668-190)/668\times100=71.6\%$. \quad
(d)~Energy saved per day $=(668-190)\times3600\times24
/1000=41\,299\,\text{Wh}=41.3\,\text{kWh}$.
Cost saved $=41.3\times150=₦6195/\text{day}$.

\textbf{B2.}
(a)~$A=4\pi r^2=4\pi\times0.01=0.1257\,\text{m}^2$.
$P_{emitted}=\sigma AT^4=5.67\times10^{-8}
\times0.1257\times500^4=5.67\times10^{-8}
\times0.1257\times6.25\times10^{10}=445\,\text{W}$.
\quad
(b)~$P_{absorbed}=\sigma AT_0^4=5.67\times10^{-8}
\times0.1257\times300^4=5.67\times10^{-8}
\times0.1257\times8.1\times10^9=57.8\,\text{W}$.
\quad
(c)~$P_{net}=445-57.8=387\,\text{W}$. \quad
(d)~$dT/dt=-P_{net}/(mc)=-387/(0.5\times800)
=-0.968°\text{C\,s}^{-1}$. \quad
(e)~As sphere cools, $T\to T_0$; $P_{net}=(T^4-T_0^4)
\to0$; deceleration decreases and sphere asymptotically
approaches ambient temperature.

\textbf{B3.}
(a)~$P_{Sun}=\sigma A_{Sun}T^4=5.67\times10^{-8}
\times4\pi(6.96\times10^8)^2\times5778^4
=3.85\times10^{26}\,\text{W}$. \quad
(b)~$I=P_{Sun}/(4\pi d^2)=3.85\times10^{26}/
(4\pi\times(1.5\times10^{11})^2)
=1361\,\text{W\,m}^{-2}$. \quad
(c)~$P_{incident}=I\times A_{panel}
=1361\times2.0=2722\,\text{W}$. \quad
(d)~$P_{absorbed}=0.92\times2722=2505\,\text{W}$.
$P_{electrical}=0.20\times2505=501\,\text{W}$. \quad
(e)~$E=P_{elec}\times t=501\times6\times3600
/1000=10.82\,\text{kWh}\approx10.8\,\text{kWh}$. \quad
(f)~Value $=10.8\times250=₦2700\,\text{/day}$.

\textbf{B4.}
(a)~$R_{glass}=0.004/(1.0\times1.2)
=3.33\times10^{-3}\,\text{K\,W}^{-1}$ (each pane).
$R_{air}=0.012/(0.025\times1.2)=0.4\,\text{K\,W}^{-1}$.
\quad
(b)~$R_{total}=2\times3.33\times10^{-3}+0.4
=0.4067\,\text{K\,W}^{-1}$. \quad
(c)~$\dot{Q}=\Delta T/R_{total}=20/0.4067
=49.2\,\text{W}$. \quad
(d)~Single glass pane:
$R_{single}=3.33\times10^{-3}\,\text{K\,W}^{-1}$.
$\dot{Q}_{single}=20/3.33\times10^{-3}=6006\,\text{W}$.
Improvement factor $=6006/49.2=122\times$. \quad
(e)~With argon ($k=0.018$):
$R_{Ar}=0.012/(0.018\times1.2)=0.556\,\text{K\,W}^{-1}$.
$R_{total,Ar}=0.5627\,\text{K\,W}^{-1}$.
$\dot{Q}_{Ar}=20/0.5627=35.5\,\text{W}$.
Better than air-filled by factor $49.2/35.5=1.38$.

\textbf{B5.}
(a)~Power absorbed by planet $=\frac{(1-\alpha)L_*
\pi R_p^2}{4\pi d^2}$.
Power emitted $=\sigma(4\pi R_p^2)T^4$.
Equating and solving:
$T=\left[\frac{(1-\alpha)L_*}{16\pi\sigma d^2}\right]^{1/4}$.
\quad
(b)~$T_{Mars}=\left[\frac{(1-0.25)\times3.85\times10^{26}}
{16\pi\times5.67\times10^{-8}\times(2.28\times10^{11})^2}\right]^{1/4}
=\left[\frac{2.89\times10^{26}}{5.92\times10^{15}}\right]^{1/4}
=(4.88\times10^{10})^{0.25}=210\,\text{K}=-63°\text{C}$.
\quad
(c)~Actual $T_{Mars}\approx210\,\text{K}$, close to our
calculation. Mars has only a thin CO$_2$ atmosphere
providing a small greenhouse effect. \quad
(d)~Increasing $\alpha$ from $0.30$ to $0.35$:
$T_{new}=T_{old}\times(1-0.35)^{0.25}/(1-0.30)^{0.25}
=T_{old}\times(0.65/0.70)^{0.25}$.
$T_{old}=255\,\text{K}$ (Earth equilibrium).
$T_{new}=255\times(0.9286)^{0.25}=255\times0.9818
=250\,\text{K}$. Temperature drops by $\sim5\,\text{K}$.
This cooling effect (increased albedo) would partially
counter global warming.

% ============================================================
%  CHAPTER 17
% ============================================================
\endgroup

\AnswersChapterBand{17}{Wave Motion}

\AnswersSectionHeading{Section A}
\begin{multicols}{2}
\begingroup
\let\textbf\AnswersTextbf
\def\quad{ --- }
\setlength{\parindent}{0pt}

\textbf{A1.~C} \quad Sound is a longitudinal mechanical
wave; oscillations are parallel to propagation.

\textbf{A2.~C} \quad $f=v/\lambda=320/0.8=400\,\text{Hz}$.

\textbf{A3.~C} \quad Maximum diffraction when
$\lambda\approx$ gap width.

\textbf{A4.~B} \quad Path difference $3\lambda/2
=(n+\tfrac{1}{2})\lambda$ with $n=1$: destructive
interference.

\textbf{A5.~C} \quad Node to adjacent antinode
$=\lambda/4$.

\textbf{A6.~C} \quad Closed pipe fundamental:
$f_{closed}=v/(4L)$. Open pipe: $f_{open}=v/(2L)
=2f_{closed}$.

\textbf{A7.~C} \quad $v=\sqrt{T/\mu}$; to double $v$,
need $T' =4T$ (quadruple tension).

\textbf{A8.~C} \quad Frequency is determined by the
source and remains unchanged when the wave crosses
a boundary between media.

\textbf{A9.~C} \quad $f_3=3v/(2L)=3\times300/(2\times1.5)
=900/3=300\,\text{Hz}$.

\textbf{A10.~B} \quad EM waves are oscillating fields,
not mechanical disturbances; they propagate through
vacuum without a medium.

\endgroup
\end{multicols}

\AnswersSectionHeading{Section B}
\begingroup
\let\textbf\AnswersTextbf
\setlength{\parindent}{0pt}
\def\quad{\par\smallskip}
% (Parenthesis activation removed: it broke ordinary parentheses in text/math.)

\textbf{B1.}
(a)~$f=1/T=1/2.0=0.5\,\text{Hz}$;
$v=f\lambda=0.5\times3.0=1.5\,\text{m\,s}^{-1}$. \quad
(b)~Path difference $=7.5/3.0=2.5\lambda$;
$2.5=(n+\tfrac{1}{2})\lambda$ with $n=2$;
points are in \textbf{antiphase}. \quad
(c)~$y=0.15\sin(2\pi ft-2\pi x/\lambda)
=0.15\sin(\pi t-2\pi x/3)\,\text{m}$. \quad
(d)~At $t=0.5\,\text{s}$, $x=0$:
$y=0.15\sin(0.5\pi)=0.15\times1=0.15\,\text{m}$.

\textbf{B2.}
(a)~$v=\sqrt{T/\mu}=\sqrt{80/0.005}=\sqrt{16000}
=126.5\,\text{m\,s}^{-1}$. \quad
(b)~$f_1=v/(2L)=126.5/(2\times1.20)=52.7\,\text{Hz}$. \quad
(c)~$f_2=2f_1=105\,\text{Hz}$; $f_3=158\,\text{Hz}$;
$f_4=211\,\text{Hz}$. \quad
(d)~For $f_1'=160\,\text{Hz}$: $v'=2L\times160
=2\times1.20\times160=384\,\text{m\,s}^{-1}$.
$T'=v'^2\times\mu=384^2\times0.005=737\,\text{N}$. \quad
(e)~Effective length $=0.90\,\text{m}$:
$f'=126.5/(2\times0.90)=70.3\,\text{Hz}$.

\textbf{B3.}
(a)~Fundamental of open pipe (flute):
$f_1=v/(2L)=340/(2\times0.60)=283\,\text{Hz}$. \quad
(b)~Harmonics: $283$, $566$, $849$, $1132\,\text{Hz}$.
\quad
(c)~Clarinet (closed pipe):
$f_1=v/(4L)=340/(4\times0.60)=141.7\,\text{Hz}$;
harmonics (odd only): $141.7$, $425$, $709\,\text{Hz}$. \quad
(d)~Clarinet lacks even harmonics; this changes the
timbre (quality). Even with the same fundamental pitch,
the waveform differs, making the instruments sound
distinct. \quad
(e)~Blowing harder produces the next odd harmonic
(overblowing); flute overblows to the octave (second
harmonic); clarinet overblows to the twelfth (third
harmonic), which is why the clarinet appears higher.

\textbf{B4.}
(a)~$\lambda=v/f=340/680=0.50\,\text{m}$. \quad
(b)~$S_2P-S_1P=4.12-4.0=0.12\,\text{m}$. \quad
(c)~$0.12/0.50=0.24\lambda$; not a whole number
or half-integer multiple. Partial (neither fully
constructive nor fully destructive) interference. \quad
(d)~$0.75/0.50=1.5\lambda=(n+\tfrac{1}{2})\lambda$;
\textbf{destructive} interference. \quad
(e)~Sketch: alternating loud (constructive) and
quiet (destructive) regions at regular spacing along
the line, with louds and quiets separated by
$\lambda/(2\sin\theta)\approx$ fixed distance depending
on geometry.

\textbf{B5.}
(a)~Over hemisphere: $I=P/(2\pi r^2)$. \quad
(b)~$I=500/(2\pi\times10^6)=7.96\times10^{-5}
\,\text{W\,m}^{-2}$. \quad
(c)~$I_{min}=10^{-6}\,\text{W\,m}^{-2}$;
$r_{max}=\sqrt{P/(2\pi I_{min})}
=\sqrt{500/(2\pi\times10^{-6})}
=\sqrt{7.96\times10^7}=8920\,\text{m}\approx8.9\,\text{km}$.
\quad
(d)~With $50\%$ absorption per $200\,\text{m}$:
$I=I_0e^{-\ln2\times r/200}\times 1/(2\pi r^2)$.
Setting $=I_{min}$ gives $r\approx1$--$2\,\text{km}$.
\quad
(e)~Sound diffracts around headlands and buildings;
the low-frequency component of the foghorn ($\sim60$--
$300\,\text{Hz}$, $\lambda\sim1$--$6\,\text{m}$)
diffracts into shadow regions unreachable by direct
sound.

% ============================================================
%  CHAPTER 18
% ============================================================
\endgroup

\AnswersChapterBand{18}{Sound Waves}

\AnswersSectionHeading{Section A}
\begin{multicols}{2}
\begingroup
\let\textbf\AnswersTextbf
\def\quad{ --- }
\setlength{\parindent}{0pt}

\textbf{A1.~C} \quad Sound travels fastest through
steel (a rigid solid with tightly packed atoms).

\textbf{A2.~C} \quad Pitch is determined by frequency.

\textbf{A3.~C} \quad $L=10\log(10^{-8}/10^{-12})
=10\log10^4=40\,\text{dB}$.

\textbf{A4.~C} \quad $f=600\times340/(340-40)
=600\times340/300=680\,\text{Hz}$.

\textbf{A5.~C} \quad $d=vt/2=340\times0.6/2
=102\,\text{m}$.

\textbf{A6.~C} \quad Ultrasound: frequency above
$20\,\text{kHz}$.

\textbf{A7.~C} \quad Source receding: wavefronts
stretched, wavelength increases, frequency decreases.

\textbf{A8.~B} \quad No molecules $\Rightarrow$ no
medium for pressure waves to propagate.

\textbf{A9.~B} \quad $70\,\text{dB}$ from $10$ sources:
$I_{total}=10^{-5}\,\text{W\,m}^{-2}$;
$I_{one}=10^{-6}\,\text{W\,m}^{-2}$.
$L=10\log(10^{-6}/10^{-12})=60\,\text{dB}$.

\textbf{A10.~D} \quad Closed pipe: first resonance
$L_1=\lambda/4$; second $L_2=3\lambda/4$.
$\lambda=2(L_2-L_1)$.

\endgroup
\end{multicols}

\AnswersSectionHeading{Section B}
\begingroup
\let\textbf\AnswersTextbf
\setlength{\parindent}{0pt}
\def\quad{\par\smallskip}
% (Parenthesis activation removed: it broke ordinary parentheses in text/math.)

\textbf{B1.}
(a)~$v=331+0.6\times25=346\,\text{m\,s}^{-1}$. \quad
(b)~Echo time $=2d/v=2\times85/346=0.491\,\text{s}$. \quad
(c)~Frequency of clapping $=1/0.491=2.04\,\text{claps\,s}^{-1}$.
\quad
(d)~New distance $=75\,\text{m}$: echo time
$=150/346=0.434\,\text{s}$; frequency
$=2.31\,\text{claps\,s}^{-1}$.

\textbf{B2.}
(a)~$f_{approach}=720\times340/(340-25)
=720\times340/315=777\,\text{Hz}$. \quad
(b)~$f_{recede}=720\times340/(340+25)
=720\times340/365=670\,\text{Hz}$. \quad
(c)~$\Delta f=777-670=107\,\text{Hz}$. \quad
(d)~Observer running away at $5\,\text{m\,s}^{-1}$,
train approaching at $25\,\text{m\,s}^{-1}$:
$f=720\times(340-5)/(340-25)=720\times335/315
=765\,\text{Hz}$. \quad
(e)~Both approaching:
$f=720\times(340+5)/(340-25)=720\times345/315
=789\,\text{Hz}$.

\textbf{B3.}
(a)~$\lambda=2(L_2-L_1)=2(56.8-18.6)\times10^{-2}
=2\times38.2\times10^{-2}=0.764\,\text{m}$. \quad
(b)~$v=f\lambda=440\times0.764=336\,\text{m\,s}^{-1}$.
\quad
(c)~End correction $e=\lambda/4-L_1
=0.191-0.186=0.005\,\text{m}=5\,\text{mm}$. \quad
(d)~$v=331+0.6\theta$;
$336=331+0.6\theta\Rightarrow\theta=8.3°\text{C}$. \quad
(e)~Third resonance: $L_3=5\lambda/4=5\times0.764/4
=0.955\,\text{m}$ (from closed end, accounting for end
correction: $L_3=5\lambda/4-e=0.950\,\text{m}$).

\textbf{B4.}
(a)~$\Delta f=24.000\,48-24.000=480\,\text{Hz}
=4.8\times10^{-4}\,\text{GHz}$. \quad
(b)~$v_{car}=c\Delta f/(2f_0)
=3\times10^8\times480/(2\times24\times10^9)
=144\times10^9/(48\times10^9)=3.0\,\text{m\,s}^{-1}$. \quad
(c)~$3.0\times3600/1000=10.8\,\text{km\,h}^{-1}$. \quad
(d)~Not speeding at $10.8\,\text{km\,h}^{-1}$. \quad
(e)~Microwaves travel much faster and in a directed beam;
sound is easily masked by traffic noise; microwave
radar is unambiguous, highly directional, works in
all weather.

\textbf{B5.}
(a)~$\lambda=v/f=1540/(5\times10^6)
=3.08\times10^{-4}\,\text{m}=0.308\,\text{mm}$. \quad
(b)~Minimum resolvable feature $\approx\lambda/2
=0.154\,\text{mm}=154\,\mu\text{m}$. \quad
(c)~$d=vt/2=1540\times80\times10^{-6}/2
=0.0616\,\text{m}=6.16\,\text{cm}$. \quad
(d)~Higher frequency $\Rightarrow$ shorter wavelength
$\Rightarrow$ better resolution (smaller features
detectable). But higher frequency also attenuates more
rapidly in tissue; the shorter wavelengths are absorbed
before reaching deep structures. Trade-off between
resolution and depth. \quad
(e)~$\Delta f/f_0=2v_{blood}/v_{sound}$;
$\Delta f=2\times0.30\times5\times10^6/1540
=1948\,\text{Hz}\approx1950\,\text{Hz}$.

% ============================================================
%  CHAPTER 19
% ============================================================
\endgroup

\AnswersChapterBand{19}{Reflection of Light}

\AnswersSectionHeading{Section A}
\begin{multicols}{2}
\begingroup
\let\textbf\AnswersTextbf
\def\quad{ --- }
\setlength{\parindent}{0pt}

\textbf{A1.~B} \quad Both angles of incidence and
reflection are measured from the normal to the
reflecting surface.

\textbf{A2.~B} \quad Plane mirror: virtual, erect,
same size, same distance behind mirror as object
is in front.

\textbf{A3.~B} \quad $1/v=1/f-1/u=1/15-1/20
=(4-3)/60=1/60$; $v=+60\,\text{cm}$.

\textbf{A4.~C} \quad Convex mirror always produces
a virtual, erect, diminished image for any object
position.

\textbf{A5.~C} \quad $f=R/2=40/2=20\,\text{cm}$.

\textbf{A6.~C} \quad Object at $F$: $1/v
=1/f-1/f=0$; $v=\infty$ (image at infinity).

\textbf{A7.~B} \quad $m=-3$: magnitude $3$ (three
times larger); negative sign means inverted.

\textbf{A8.~C} \quad $n=360°/60°-1=6-1=5$ images.

\textbf{A9.~B} \quad $u=6\,\text{cm}<f=12\,\text{cm}$;
object inside focal length.
$1/v=1/12-1/6=1/12-2/12=-1/12\Rightarrow v=-12\,\text{cm}$
(virtual, $12\,\text{cm}$ behind mirror).
$m=-v/u=-(-12)/6=+2$; erect and magnified.
Answer: virtual, erect, at $12\,\text{cm}$ behind mirror.

\textbf{A10.~C} \quad Convex mirror: wide field of view
and always produces an upright image regardless of
object distance.

\endgroup
\end{multicols}

\AnswersSectionHeading{Section B}
\begingroup
\let\textbf\AnswersTextbf
\setlength{\parindent}{0pt}
\def\quad{\par\smallskip}
% (Parenthesis activation removed: it broke ordinary parentheses in text/math.)

\textbf{B1.}
$f=20\,\text{cm}$ (concave).
(a)~$u=60$: $1/v=1/20-1/60=3/60-1/60=2/60$;
$v=30\,\text{cm}$; $m=-1/2$. Real, inverted,
diminished. \quad
(b)~$u=40$: $1/v=1/20-1/40=1/40$; $v=40\,\text{cm}$;
$m=-1$. Real, inverted, same size. \quad
(c)~$u=30$: $1/v=1/20-1/30=3/60-2/60=1/60$;
$v=60\,\text{cm}$; $m=-2$. Real, inverted, magnified. \quad
(d)~$u=10$: $1/v=1/20-1/10=1/20-2/20=-1/20$;
$v=-20\,\text{cm}$; $m=+2$. Virtual, erect, magnified,
$20\,\text{cm}$ behind mirror.

\textbf{B2.}
(a)~$1/v=1/f-1/u=1/(-30)-1/200=-1/30-1/200$.
LCM 600: $=-20/600-3/600=-23/600$;
$v=-600/23=-26.1\,\text{cm}$ (virtual). \quad
(b)~$m=+v/u=26.1/200=+0.130$. \quad
(c)~Image height $=0.130\times1.8=0.234\,\text{m}
=23.4\,\text{cm}$. \quad
(d)~The diminished image looks like a distant object;
our brain judges distance by apparent size, so objects
in convex mirrors appear farther away than they are ---
hence the warning ``Objects in mirror are closer than
they appear.''

\textbf{B3.}
(a)~For $m=+3$ (virtual, erect): $m=-v/u$
(real-is-positive for mirror) $\Rightarrow v=-3u$.
$1/f=1/u+1/v=1/u-1/(3u)=2/(3u)$;
$u=2f/3=2\times2.4/3=1.6\,\text{cm}$. \quad
(b)~$v=-3\times1.6=-4.8\,\text{cm}$ (virtual,
$4.8\,\text{cm}$ behind mirror). \checkmark \quad
(c)~Virtual image is upright and magnified; the
dentist sees the tooth right way up through the mirror,
larger than actual. A real image would be on the same
side as the patient and inverted --- useless. \quad
(d)~$u=3.0\,\text{cm}$, $f=2.4\,\text{cm}$.
$1/v=1/2.4-1/3=5/12-4/12=1/12$; $v=12\,\text{cm}$.
Real, inverted, magnified ($m=-4$). Not suitable ---
image is inverted and far from mirror.

\textbf{B4.}
(a)~Object at $F$: image at infinity (parallel beam).
\quad
(b)~As $u\to f$: $1/v=1/f-1/u\to0$; $v\to\infty$. \quad
(c)~$R=2f=2\times3.5=7.0\,\text{cm}$. \quad
(d)~If filament is $0.5\,\text{cm}$ beyond $F$:
$u=4.0\,\text{cm}$; $v=1/f-1/u=1/3.5-1/4
=0.0357$ $\Rightarrow v=28\,\text{cm}$ (real, inverted
image at $28\,\text{cm}$ in front of mirror);
the beam is convergent, not parallel --- beam narrows
to a spot then diverges. \quad
(e)~Spherical reflector has spherical aberration;
parallel rays further from axis focus at a different
point than paraxial rays (the caustic). A parabolic
reflector focuses all parallel rays to a single perfect
focus.

\textbf{B5.}
(a)~$1/u$ values: $0.0667$, $0.0500$, $0.0400$,
$0.0333$, $0.0250$.
$1/v$ values: $0.0167$, $0.0333$, $0.0467$,
$0.0541$, $0.0649$. \quad
(b)--(d)~Intercepts on $1/v$ and $1/u$ axes both
$\approx1/f$. Plot should give gradient $\approx-1$
and intercepts $\approx1/f$.
Average: $f\approx12.0$--$12.5\,\text{cm}$. \quad
(e)~Sources of error: parallax in reading scale,
locating exact image position by screen, spherical
aberration for rays far from axis.

% ============================================================
%  CHAPTER 20
% ============================================================
\endgroup

\AnswersChapterBand{20}{Refraction and Optical Instruments}

\AnswersSectionHeading{Section A}
\begin{multicols}{2}
\begingroup
\let\textbf\AnswersTextbf
\def\quad{ --- }
\setlength{\parindent}{0pt}

\textbf{A1.~B} \quad Light going from less dense
($n=1.33$) to more dense ($n=1.50$) bends toward
the normal.

\textbf{A2.~C} \quad $\sin\theta_c=1/n$;
$n=1/\sin42°=1/0.6691=1.494\approx1.49$.

\textbf{A3.~C} \quad Converging lens, $u=15$, $f=10$:
$1/v=1/10-1/15=1/30$; $v=+30\,\text{cm}$.

\textbf{A4.~B} \quad Diverging lens always produces
virtual, erect, diminished image.

\textbf{A5.~C} \quad $P=1/f=1/0.20=5.0\,\text{D}$.

\textbf{A6.~B} \quad TIR conditions: denser to less
dense medium; angle of incidence $>$ critical angle.

\textbf{A7.~B} \quad $M=f_o/f_e$: increase $f_o$
and decrease $f_e$.

\textbf{A8.~C} \quad Short-sight corrected by diverging
lens (brings image forward onto retina).

\textbf{A9.~B} \quad $\lambda_{medium}=\lambda_{air}/n$;
decreases by factor $n$.

\textbf{A10.~B} \quad Both lenses short focal length:
$M_{obj}=v_o/u_o$ (large when $u_o\approx f_o$);
$M_{eye}=D/f_e$ (large when $f_e$ small).

\endgroup
\end{multicols}

\AnswersSectionHeading{Section B}
\begingroup
\let\textbf\AnswersTextbf
\setlength{\parindent}{0pt}
\def\quad{\par\smallskip}
% (Parenthesis activation removed: it broke ordinary parentheses in text/math.)

\textbf{B1.}
(a)~$\sin\theta_c=n_{liquid}/n_{glass}=1.40/1.60=0.875$;
$\theta_c=61.0°$. \quad
(b)~$50°<61°$; TIR does \textit{not} occur. Ray
refracts into liquid. \quad
(c)~$\sin\theta_r=(1.60/1.40)\sin50°=1.143\times0.766
=0.875\Rightarrow\theta_r=61.0°$. \quad
(d)~Glass-air TIR: $\sin\theta_c=1/1.60=0.625$;
$\theta_c=38.7°$. The ray must strike the glass-air
surface at $>38.7°$ from the normal, which means it
travels at $<90°-38.7°=51.3°$ to the surface.

\textbf{B2.}
(a)~For $m=5$: $m=v/u\Rightarrow v=5u$.
$1/f=1/v-1/u$ (Cartesian, $u$ negative for real object)
$\Rightarrow1/15=1/(5u)+1/u=6/(5u)\Rightarrow u=18\,\text{cm}$.
Image distance $=5\times18=90\,\text{cm}$. \quad
(b)~Object at $18\,\text{cm}$. \quad
(c)~Screen at $90\,\text{cm}$ from lens. \quad
(d)~$M=1+D/f=1+25/15=2.67\times$. \quad
(e)~Projection: image $5\times3=15\,\text{cm}$ tall on
screen; magnifying glass: virtual image at near point,
$M\approx2.67$, image height $=2.67\times3=8\,\text{cm}$.
Projection gives larger absolute image height.

\textbf{B3.}
(a)~$u_o=1.0\,\text{cm}$, $f_o=0.80\,\text{cm}$:
$1/v_o=1/0.80-1/1.0=1.25-1.0=0.25$;
$v_o=4.0\,\text{cm}$. \quad
(b)~$m_o=v_o/u_o=4.0/1.0=4.0\times$. \quad
(c)~Eyepiece (normal adjustment): $M_e=D/f_e
=25/3.0=8.33\times$. \quad
(d)~$M_{total}=4.0\times8.33=33.3\times$. \quad
(e)~Microscope objective: short $f_o$ because the
object is placed just beyond $f_o$ to produce a
\textit{large} real image inside the tube (large $v_o$
means large $m_o$). Telescope objective: long $f_o$
because the object is at infinity; a long $f_o$
produces a large focal-plane image of a distant object,
giving high angular resolution and easier eyepiece
magnification.

\textbf{B4.}
(a)~Long-sighted, near point $50\,\text{cm}$:
$u=-25\,\text{cm}$ (object at reading distance),
$v=-50\,\text{cm}$ (virtual image at near point).
$1/f=1/v-1/u=1/(-50)-1/(-25)=-0.02+0.04=+0.02$;
$P=+2.0\,\text{D}$. \quad
(b)~Short-sighted, far point $2.0\,\text{m}$:
$u=\infty$, $v=-2.0\,\text{m}$;
$P=1/f=1/v=1/(-2.0)=-0.5\,\text{D}$. \quad
(c)~Combined: $P=+2.5+(-1.5)=+1.0\,\text{D}$. \quad
(d)~Contact lenses sit on the cornea; the distances to
the objects and images change because the lens is at a
different position. The effective object and image
distances are measured from the cornea, not from a
spectacle position $\sim1.5\,\text{cm}$ in front.
This leads to a different required power (typically
slightly less power for myopia correction and slightly
more for hypermetropia).

\textbf{B5.}
(a)~$\sin\theta_c=1.45/1.55=0.9355$;
$\theta_c=69.4°$. \quad
(b)~In core: $v_1=c/n_1=3\times10^8/1.55
=1.935\times10^8\,\text{m\,s}^{-1}$.
In cladding: $v_2=3\times10^8/1.45
=2.069\times10^8\,\text{m\,s}^{-1}$. \quad
(c)~$\lambda_{core}=\lambda_{vac}/n_1
=1550/1.55=1000\,\text{nm}$. \quad
(d)~Every $15\,\text{km}$: intensity halved.
After $100\,\text{km}$: $n=100/15=6.67$ half-lives.
$I/I_0=(0.5)^{6.67}=0.5^{6.67}=0.010$.
In dB: $10\log(0.010)=-20\,\text{dB}$. \quad
(e)~Each amplifier restores signal. Every $80\,\text{km}$:
$80/15=5.33$ half-lives; $I/I_0=0.5^{5.33}=0.0247
=2.47\%$ before each amplifier. After amplification,
full power restored. For $1000\,\text{km}$:
$1000/80=12.5$, so $12$ amplifiers needed (one every
$80\,\text{km}$ for segments $80, 160,...,960$; the
last $40\,\text{km}$ brings intensity to
$0.5^{40/15}\times100\%=15.9\%>1\%$, just adequate).

% ============================================================
%  CHAPTER 21
% ============================================================
\endgroup

\AnswersChapterBand{21}{Electrostatics}

\AnswersSectionHeading{Section A}
\begin{multicols}{2}
\begingroup
\let\textbf\AnswersTextbf
\def\quad{ --- }
\setlength{\parindent}{0pt}

\textbf{A1.~C} \quad Original $F=k(2Q)(3Q)/d^2
=6kQ^2/d^2$. New: charges $4Q,6Q$; distance $d/2$:
$F'=k(4Q)(6Q)/(d/2)^2=96kQ^2/d^2=16F$.

\textbf{A2.~B} \quad Inside a conductor in
electrostatic equilibrium: $E=0$ throughout.

\textbf{A3.~B} \quad Work done by field on charge $q$
moving from $V_1$ to $V_2$:
$W=q(V_1-V_2)$ (field does work moving charge
from high to low potential).

\textbf{A4.~C} \quad $E=V/d=250/(5\times10^{-3})
=5\times10^4\,\text{V\,m}^{-1}$.

\textbf{A5.~B} \quad $V=kQ/r\propto1/r$.

\textbf{A6.~B} \quad Positive rod induces negative
charges to move toward cap; leaves lose some negative
charge (become less negative); leaves diverge less
$\Rightarrow$ come together.

\textbf{A7.~B} \quad Charging by induction produces
charge opposite to the inducing charge.

\textbf{A8.~B} \quad $1\,\text{V\,m}^{-1}
=1\,\text{J\,C}^{-1}\text{m}^{-1}=1\,\text{N\,C}^{-1}$.
The unit is $\text{V\,m}^{-1}$.

\textbf{A9.~C} \quad $F_E/F_G\approx10^{39}$ for
electron-proton pair.

\textbf{A10.~B} \quad Equipotentials are always
perpendicular to field lines; no work done moving
along an equipotential.

\endgroup
\end{multicols}

\AnswersSectionHeading{Section B}
\begingroup
\let\textbf\AnswersTextbf
\setlength{\parindent}{0pt}
\def\quad{\par\smallskip}
\catcode`(=\active
\def({\AnswersParen}

\textbf{B1.}
(a)~$F_{12}=kQ_1Q_2/r_{12}^2
=9\times10^9\times4\times10^{-6}\times2\times10^{-6}
/0.09=0.8\,\text{N}$ (attractive, toward $+x$). \quad
(b)~$F_{32}=9\times10^9\times3\times10^{-6}
\times2\times10^{-6}/0.09=0.6\,\text{N}$ (attractive,
toward $-x$). \quad
(c)~Net force on $Q_2=0.8-0.6=0.2\,\text{N}$
toward $Q_1$ (in $-x$ direction). \quad
(d)~Zero field point between $Q_1$ and $Q_3$: let
distance from $Q_1$ be $x$.
$kQ_1/x^2=kQ_3/(0.6-x)^2$;
$4/(x)^2=3/(0.6-x)^2$;
$2/x=\sqrt{3}/(0.6-x)$;
$2(0.6-x)=x\sqrt{3}$;
$1.2-2x=1.732x$;
$x=1.2/3.732=0.321\,\text{m}$ from $Q_1$.

\textbf{B2.}
(a)~$E=V/d=400/(2\times10^{-3})
=2\times10^5\,\text{V\,m}^{-1}$. \quad
(b)~$\sigma=\varepsilon_0E=8.85\times10^{-12}
\times2\times10^5=1.77\times10^{-6}\,\text{C\,m}^{-2}
=1.77\,\mu\text{C\,m}^{-2}$. \quad
(c)~$Q=\sigma A=1.77\times10^{-6}\times0.050
=8.85\times10^{-8}\,\text{C}=88.5\,\text{nC}$. \quad
(d)~$a=eE/m=1.6\times10^{-19}\times2\times10^5/
(9.11\times10^{-31})=3.51\times10^{16}\,\text{m\,s}^{-2}$.
$v^2=2\times3.51\times10^{16}\times2\times10^{-3}$
$=1.404\times10^{14}$; $v=1.19\times10^7\,\text{m\,s}^{-1}$.

\textbf{B3.}
(a)~$E_{surface}=kQ/R^2
=9\times10^9\times5\times10^{-6}/(0.10)^2
=4.5\times10^6\,\text{V\,m}^{-1}$. \quad
(b)~$E_{0.30\,m}=9\times10^9\times5\times10^{-6}/0.09
=5\times10^5\,\text{V\,m}^{-1}$. \quad
(c)~$V=kQ/R=9\times10^9\times5\times10^{-6}/0.10
=4.5\times10^5\,\text{V}$. \quad
(d)~$W=q\Delta V=2\times10^{-6}\times4.5\times10^5
=0.90\,\text{J}$. \quad
(e)~By Gauss's Law: a Gaussian surface inside the
spherical shell encloses zero charge
($Q_{enc}=0\Rightarrow E=0$).

\textbf{B4.}
(a)~Forces: weight $mg$ downward;
electric force $qE$ upward. \quad
(b)~For equilibrium: $qE=mg$. \quad
(c)~$E=V/d=490/(6\times10^{-3})
=8.167\times10^4\,\text{V\,m}^{-1}$.
$q=mg/E=3.27\times10^{-15}\times10/8.167\times10^4
=4.0\times10^{-19}\,\text{C}$. \quad
(d)~$n=q/e=4.0\times10^{-19}/1.6\times10^{-19}
=2.5$... Not an integer. Closest: $3e$ or $2e$.
Rounding: $q=4.0\times10^{-19}\approx2.5e$;
likely $q=3e$ with slight experimental error. \quad
(e)~Electric force doubles; net upward force acts;
drop accelerates upward.

\textbf{B5.}
(a)~(i)~At sharp points, charge density and field are
highest; air ionises more readily. (ii)~Ionised air
forms a conductive path (corona discharge). (iii)~When
lightning strikes the rod, current flows down the
copper conductor to earth, bypassing the building. \quad
(b)~$I=Q/t=5.0/0.001=5000\,\text{A}$. \quad
(c)~$E=QV=5.0\times3.0\times10^8
=1.5\times10^9\,\text{J}$. \quad
(d)~$E=1.5\times10^9/(3.6\times10^6)
=417\,\text{kWh}$. At $60\,\text{W}$:
$t=1.5\times10^9/(60)=2.5\times10^7\,\text{s}
\approx289\,\text{days}$. \quad
(e)~Lightning strikes are brief, unpredictable,
and spatially random; the energy per strike is
enormous but arrives in $\sim1\,\text{ms}$; storage
systems cannot handle such instantaneous power;
harvesting the energy economically is not feasible.

% ============================================================
%  CHAPTER 22
% ============================================================
\endgroup

\AnswersChapterBand{22}{Capacitance}

\AnswersSectionHeading{Section A}
\begin{multicols}{2}
\begingroup
\let\textbf\AnswersTextbf
\def\quad{ --- }
\setlength{\parindent}{0pt}

\textbf{A1.~A} \quad $C=Q/V=60\times10^{-6}/12
=5.0\,\mu\text{F}$.

\textbf{A2.~D} \quad Three in series:
$1/C=3/6=0.5$; $C=2\,\mu\text{F}$.

\textbf{A3.~B} \quad $E=\tfrac{1}{2}CV^2
=\tfrac{1}{2}\times50\times10^{-6}\times(200)^2
=1.0\,\text{J}$.

\textbf{A4.~C} \quad Battery disconnected: $Q$ constant.
Inserting dielectric increases $C$; $V=Q/C$
decreases.

\textbf{A5.~A} \quad $Q=C_1V_1=10\times5=50\,\mu\text{C}$.
$V_f=Q/(C_1+C_2)=50/50=1.0\,\text{V}$.

\textbf{A6.~C} \quad $I=V/R=(V_0e^{-t/RC})/R$;
decreases exponentially.

\textbf{A7.~B} \quad $\tau=RC=10\times10^3\times470
\times10^{-6}=4.7\,\text{s}$.

\textbf{A8.~B} \quad $C=\varepsilon A/d$.
Double $A$, halve $d$: $C'=\varepsilon(2A)/(d/2)
=4\varepsilon A/d=4C$ (quadruples).

\textbf{A9.~C} \quad $e^{-3}\approx0.050=5\%$
of initial voltage.

\textbf{A10.~C} \quad Total $C=2C$;
$E=\tfrac{1}{2}(2C)V^2=CV^2
=2\times(\tfrac{1}{2}CV^2)$. Double.

\endgroup
\end{multicols}

\AnswersSectionHeading{Section B}
\begingroup
\let\textbf\AnswersTextbf
\setlength{\parindent}{0pt}
\def\quad{\par\smallskip}
\catcode`(=\active
\def({\AnswersParen}

\textbf{B1.}
(a)~$A=200\,\text{cm}^2=0.02\,\text{m}^2$;
$C=\varepsilon_0A/d=8.85\times10^{-12}
\times0.02/(1.5\times10^{-3})=118\,\text{pF}$. \quad
(b)~$Q=CV=118\times10^{-12}\times600=70.8\,\text{nC}$;
$E=\tfrac{1}{2}CV^2=\tfrac{1}{2}\times118\times10^{-12}
\times3.6\times10^5=2.12\times10^{-5}\,\text{J}
=21.2\,\mu\text{J}$. \quad
(c)~With dielectric: $C'=4.5\times118=531\,\text{pF}$.
$Q'=531\times10^{-12}\times600=318.6\,\text{nC}$;
$E'=\tfrac{1}{2}\times531\times10^{-12}\times600^2
=95.6\,\mu\text{J}$. \quad
(d)~Battery maintains constant $V$; more charge is
drawn from battery as $C$ increases; extra energy
$=95.6-21.2=74.4\,\mu\text{J}$ comes from the battery.
Half of this extra energy goes into polarising the
dielectric molecules (dielectric work); half stores in
field. The battery supplies all the extra energy.

\textbf{B2.}
(a)~$C_{12}=C_1+C_2=3+6=9\,\mu\text{F}$. \quad
(b)~$C_{123}$: $9\,\mu\text{F}$ in series with
$4\,\mu\text{F}$:
$1/C=1/9+1/4=13/36$; $C_{123}=36/13=2.77\,\mu\text{F}$.
\quad
(c)~$C_{total}=C_{123}+C_4=2.77+8=10.77\,\mu\text{F}$.
\quad
(d)~$Q_{total}=C_{total}\times V=10.77\times10^{-6}
\times20=215.4\,\mu\text{C}$. \quad
(e)~$Q$ through $C_4=C_4\times V=8\times20
=160\,\mu\text{C}$.
$Q$ through $C_{123}=215.4-160=55.4\,\mu\text{C}$.
$V_{C3}=Q_{123}/C_{123}\times C_{123}/C_3$...
Actually $Q_{series}=Q_{12series with 3}=55.4\,\mu\text{C}$.
$V_{C3}=55.4/4=13.85\,\text{V}$.

\textbf{B3.}
(a)~$Q_0=CV_0=1000\times10^{-6}\times15
=0.015\,\text{C}$;
$I_0=V_0/R=15/(5000)=3\,\text{mA}$;
$E_0=\tfrac{1}{2}CV_0^2=\tfrac{1}{2}\times10^{-3}
\times225=0.1125\,\text{J}$. \quad
(b)~$\tau=RC=5000\times10^{-3}=5\,\text{s}$. \quad
(c)~$5=15e^{-t/5}\Rightarrow e^{-t/5}=1/3
\Rightarrow t=5\ln3=5.49\,\text{s}$. \quad
(d)~$I=I_0e^{-t/\tau}=3\times10^{-3}e^{-10/5}
=3\times10^{-3}\times e^{-2}=3\times0.135
=0.406\,\text{mA}$. \quad
(f)~(i)~$I_0=V_0/R'=15/10000=1.5\,\text{mA}$ (halved).
(ii)~$\tau'=10000\times10^{-3}=10\,\text{s}$ (doubled).
(iii)~Total charge $=Q_0=CV_0=0.015\,\text{C}$
(unchanged; all charge eventually delivered regardless
of $R$).

\textbf{B4.}
(a)~$C=2E/V^2=2\times360/25\times10^6
=28.8\,\mu\text{F}$. \quad
(b)~$Q=CV=28.8\times10^{-6}\times5000
=0.144\,\text{C}$. \quad
(c)~$\tau=RC=50\times28.8\times10^{-6}
=1.44\times10^{-3}\,\text{s}\neq5\,\text{ms}$.
Discrepancy: the $\tau=5\,\text{ms}$ given implies
$C=\tau/R=5\times10^{-3}/50=100\,\mu\text{F}$.
These are different specifications; likely the
question intends two separate design scenarios. \quad
(d)~Peak current $=V/R=5000/50=100\,\text{A}$. \quad
(e)~Peak power $=V^2/R=25\times10^6/50
=500\,000\,\text{W}=500\,\text{kW}$. \quad
(f)~After $5\tau$: $99.3\%$ of energy delivered.
The final $0.7\%$ remains; for medical purposes,
essentially complete discharge in $5\tau$.

\textbf{B5.}
(a)~$\ln V$ values:
$t=0$: $\ln8.00=2.079$;
$5$: $\ln5.90=1.775$;
$10$: $1.470$; $15$: $1.163$;
$20$: $0.858$; $25$: $0.554$. \quad
(b)--(c)~Gradient $\approx-(2.079-0.554)/25
=-1.525/25=-0.061\,\text{s}^{-1}$;
$\tau=1/0.061=16.4\,\text{s}$. \quad
(d)~$y$-intercept $\approx2.08$;
$V_0=e^{2.08}=8.0\,\text{V}$. \quad
(e)~$C=\tau/R=16.4/(33\times10^3)=4.97\times10^{-4}
\,\text{F}\approx497\,\mu\text{F}$.
Close to labelled $470\,\mu\text{F}$; discrepancy
due to measurement uncertainty and tolerance
($\pm20\%$ is typical for electrolytic capacitors).

% ============================================================
%  CHAPTER 23
% ============================================================
\endgroup

\AnswersChapterBand{23}{Current Electricity}

\AnswersSectionHeading{Section A}
\begin{multicols}{2}
\begingroup
\let\textbf\AnswersTextbf
\def\quad{ --- }
\setlength{\parindent}{0pt}

\textbf{A1.~A} \quad $I=Q/t=120/(4\times60)
=120/240=0.5\,\text{A}$.

\textbf{A2.~B} \quad Parallel pair: $6\parallel6=3\,\Omega$.
Series with third: $3+6=9\,\Omega$.

\textbf{A3.~A} \quad $V_T=\mathcal{E}-Ir=6-2\times1
=4.0\,\text{V}$.

\textbf{A4.~B} \quad KVL: sum of p.d.s around a loop
$=0$; consequence of conservation of energy (path
independence of potential).

\textbf{A5.~C} \quad $V_{out}=12\times9/(3+9)
=12\times0.75=9\,\text{V}$.

\textbf{A6.~C} \quad Metal: increasing temperature
increases lattice vibration; more electron-ion
collisions; higher resistance.

\textbf{A7.~A} \quad $S=QR/P=5\times8/10=4\,\Omega$.

\textbf{A8.~C} \quad NTC thermistor: resistance
decreases with temperature (more carriers freed
thermally).

\textbf{A9.~A} \quad $v_d=I/(nqA)
=3/(8.5\times10^{28}\times1.6\times10^{-19}
\times2\times10^{-6})$
$=3/(2.72\times10^4)=1.1\times10^{-4}\,\text{m\,s}^{-1}$.

\textbf{A10.~B} \quad At balance, zero current from
cell being measured; terminal voltage $=$ EMF.

\endgroup
\end{multicols}

\AnswersSectionHeading{Section B}
\begingroup
\let\textbf\AnswersTextbf
\setlength{\parindent}{0pt}
\def\quad{\par\smallskip}
\catcode`(=\active
\def({\AnswersParen}

\textbf{B1.}
(a)~$A=\pi(0.5\times10^{-3})^2=7.85\times10^{-7}\,\text{m}^2$;
$R=\rho L/A=1.72\times10^{-8}\times2/(7.85\times10^{-7})
=0.0438\,\Omega$. \quad
(b)~$I=V/R=3.0/0.0438=68.5\,\text{A}$. \quad
(c)~$v_d=I/(nqA)=68.5/(8.5\times10^{28}
\times1.6\times10^{-19}\times7.85\times10^{-7})
=68.5/(1.066\times10^4)=6.43\times10^{-3}\,\text{m\,s}^{-1}$.
\quad
(d)~The electric field propagates at close to $c$;
all electrons start drifting almost simultaneously.
Drift is slow ($6\,\text{mm\,s}^{-1}$) but the signal
(field) travels at nearly $3\times10^8\,\text{m\,s}^{-1}$.
\quad
(e)~Half diameter $\Rightarrow$ area $\times\frac{1}{4}$;
$R\times4=0.175\,\Omega$; $I=3.0/0.175=17.1\,\text{A}$.

\textbf{B2.}
(a)~Branch currents $I_1$ (through $R_1$), $I_2$
(through $\mathcal{E}_2,r_2,R_2$), $I_3$ (through
$R_3$). KCL at junction: $I_1=I_2+I_3$. \quad
(b)~Loop 1 ($\mathcal{E}_1$, $r_1$, $R_1$, and
back through junction):
$\mathcal{E}_1=I_1(r_1+R_1)=I_1(0.5+3)=3.5I_1$.
Loop 2 ($\mathcal{E}_2$, $r_2$, $R_2$, same V as
branch 1):
$V_{R_1}=I_1R_1=I_2(R_2+r_2)+\mathcal{E}_2$. \quad
(c)~From loop 1: $I_1=10/3.5=2.857\,\text{A}$.
$V_{R_1}=2.857\times3=8.571\,\text{V}$.
$I_2=(8.571-4)/(2+1)=4.571/3=1.524\,\text{A}$.
$I_3=8.571/6=1.429\,\text{A}$.
Check: $I_2+I_3=2.953\approx I_1=2.857$. Small
discrepancy from rounding; circuit consistent. \quad
(d)~Power from $\mathcal{E}_1$: $10\times2.857
=28.57\,\text{W}$.
Power from $\mathcal{E}_2$: $4\times1.524
=6.10\,\text{W}$.
Total: $34.67\,\text{W}$.
Resistors: $I_1^2\times3.5+I_2^2\times3+I_3^2\times6
=28.57+6.97+12.24=47.78$. Discrepancy from rounding
in $I$ values; consistent at 3 significant figures.

\textbf{B3.}
(a)~$R_{23}=6\times3/(6+3)=2\,\Omega$.
$R_{total}=r+R_1+R_{23}+R_4=0.5+2+2+1=5.5\,\Omega$.
\quad
(b)~$I=12/5.5=2.182\,\text{A}$. \quad
(c)~$V_T=\mathcal{E}-Ir=12-2.182\times0.5
=10.91\,\text{V}$. \quad
(d)~$V_{23}=I\times R_{23}=2.182\times2
=4.36\,\text{V}$. \quad
(e)~$I_3=V_{23}/R_3=4.36/3=1.45\,\text{A}$. \quad
(f)~$P_{R4}=I^2R_4=2.182^2\times1=4.76\,\text{W}$.

\textbf{B4.}
(a)~Balance: $P/Q=R/S$; $R=PS/Q=100\times85/200
=42.5\,\Omega$. \quad
(b)~New $S=86.5\,\Omega$:
$R=100\times86.5/200=43.25\,\Omega$. \quad
(c)~$GF=(\Delta R/R)/\varepsilon
=0.0176/(8.8\times10^{-5})=200$.
Typical strain gauges have $GF\approx2$; this
is very high, suggesting a semiconductor strain
gauge. \quad
(d)~A bridge produces a \textit{difference} signal;
at balance output is zero. A small $\Delta R$ produces
a large unbalanced voltage proportional to $\Delta R$.
A simple divider would require measuring a small
change on a large background voltage, much less
sensitive.

\textbf{B5.}
(a)~Potential gradient $=\mathcal{E}_d/L
=4.0/100=0.04\,\text{V\,cm}^{-1}$. \quad
(b)~$\mathcal{E}_2=\text{gradient}\times\ell_2
=0.04\times45=1.80\,\text{V}$. \quad
(c)~$V_{T2}=0.04\times35=1.40\,\text{V}$ (terminal
voltage under load). \quad
(d)~$I=(\mathcal{E}_2-V_{T2})/R_{internal}$;
but we need $I$: from $V=\mathcal{E}-Ir$,
$I=(1.80-1.40)/r$. Also $V=IR_L\Rightarrow
I=V/R_L=1.40/5=0.28\,\text{A}$.
$r=(1.80-1.40)/0.28=0.40/0.28=1.43\,\Omega$. \quad
(e)~Potentiometer draws zero current at null point;
the EMF measured equals the true EMF with no $Ir$
drop. A voltmeter draws current and reads terminal
voltage $\mathcal{E}-Ir<\mathcal{E}$.

% ============================================================
%  CHAPTER 24
% ============================================================
\endgroup

\AnswersChapterBand{24}{Electrical Energy and Power}

\AnswersSectionHeading{Section A}
\begin{multicols}{2}
\begingroup
\let\textbf\AnswersTextbf
\def\quad{ --- }
\setlength{\parindent}{0pt}

\textbf{A1.~C} \quad $P=I^2R=4\times100=400\,\text{W}$.

\textbf{A2.~C} \quad $1\,\text{kWh}=1000\times3600
=3.6\times10^6\,\text{J}$.

\textbf{A3.~B} \quad $N_s>N_p\Rightarrow V_s>V_p$:
step-up transformer.

\textbf{A4.~B} \quad $P_{loss}=I^2R$; transmit at high
$V$ (low $I$) to minimise $I^2R$ losses.

\textbf{A5.~B} \quad $V_{rms}=V_0/\sqrt{2}
=100/1.414=70.7\,\text{V}$.

\textbf{A6.~B} \quad Step-down $240\to12$:
$N_p/N_s=20$; $I_s/I_p=N_p/N_s=20$;
$I_p/I_s=1/20$, i.e.\ $1:20$.

\textbf{A7.~C} \quad $E=2\times3\times30
=180\,\text{kWh}$.

\textbf{A8.~C} \quad Fuse: wire melts when current
exceeds rated value.

\textbf{A9.~C} \quad $I=P/V$; double $V\Rightarrow I$
halves; $P_{loss}=I^2R$ is quartered.

\textbf{A10.~C} \quad Transformers need changing
flux; DC gives constant flux; no induction in
secondary.

\endgroup
\end{multicols}

\AnswersSectionHeading{Section B}
\begingroup
\let\textbf\AnswersTextbf
\setlength{\parindent}{0pt}
\def\quad{\par\smallskip}
\catcode`(=\active
\def({\AnswersParen}

\textbf{B1.}
(a)~Daily kWh:
AC: $1.8\times8=14.4$; Fridge: $0.15\times24=3.6$;
Bulbs: $5\times0.015\times6=0.45$; TV: $0.12\times5
=0.60$; Laptop: $0.065\times4=0.26$; Kettle:
$1.5\times0.5=0.75$; Total $=20.06\,\text{kWh}$.
\quad
(b)~$20.06\,\text{kWh/day}$. \quad
(c)~Monthly: $20.06\times30\times80=₦48\,144$. \quad
(d)~Inverter AC: $1.2\times8=9.6\,\text{kWh}$.
Saving: $(14.4-9.6)\times30\times80=₦11\,520/\text{month}$.
\quad
(e)~Payback $=180\,000/11\,520\approx15.6$ months.

\textbf{B2.}
(a)~$P_{in}=V\times I=220\times10^3\times500
=110\,\text{MW}$. \quad
(b)~$P_{loss}=I^2R=500^2\times60=15\,\text{MW}$.
\quad
(c)~$\Delta V=IR=500\times60=30\,000\,\text{V}
=30\,\text{kV}$. \quad
(d)~$V_{Lagos}=220-30=190\,\text{kV}$. \quad
(e)~$P_L=190\times10^3\times500=95\,\text{MW}$. \quad
(f)~$P_{in}=P_{loss}+P_L=15+95=110\,\text{MW}$. \checkmark

\textbf{B3.}
(a)~$N_p/N_s=V_p/V_s=240/12=20$;
$N_p=20\times60=1200\,\text{turns}$. \quad
(b)~$I_p=I_sN_s/N_p=5\times60/1200=0.25\,\text{A}$.
\quad
(c)~$P_{in}=P_{out}=V_sI_s=12\times5=60\,\text{W}$.
\quad
(d)~$\eta=92\%$: $P_{in}=60/0.92=65.2\,\text{W}$;
$I_p=65.2/240=0.272\,\text{A}$. \quad
(e)~Heat per hour $=(65.2-60)\times3600
=18\,720\,\text{J}$. \quad
(f)~Losses: (1) Eddy current losses in iron core
(reduced by lamination); (2) Hysteresis losses (reduced
by using soft iron); (3) Copper losses ($I^2R$ in
windings); (4) Flux leakage.

\textbf{B4.}
(a)~$v(t)=325\sin(100\pi t)\,\text{V}$. \quad
(b)~$V_{rms}=325/\sqrt{2}=229.8\approx230\,\text{V}$.
\quad
(c)~$I_0=V_0/R=325/40=8.125\,\text{A}$;
$I_{rms}=8.125/\sqrt{2}=5.74\,\text{A}$;
$P=V_{rms}I_{rms}=230\times5.74=1320\,\text{W}$. \quad
(d)~Voltmeter reads $V_{rms}=230\,\text{V}$;
ammeter reads $I_{rms}=5.74\,\text{A}$. \quad
(e)~$V_{s,rms}=V_{p,rms}\times N_s/N_p
=230\times50/200=57.5\,\text{V}$;
$V_{s,peak}=57.5\times\sqrt{2}=81.3\,\text{V}$.

\textbf{B5.}
(a)~Max power $=VI=240\times13=3120\,\text{W}$. \quad
(b)~$I_{total}=(1800+900+1200+60)/240
=3960/240=16.5\,\text{A}$. \quad
(c)~$16.5>13\,\text{A}$: fuse \textit{will} blow.
Fitting a $20\,\text{A}$ fuse is dangerous: fuse no
longer protects the cable (rated $13\,\text{A}$);
cables may overheat and catch fire. \quad
(d)~Total current through cord $=16.5\,\text{A}$;
$P_{cord}=I^2R_{cord}=16.5^2\times0.5
=136\,\text{W}$. \quad
(e)~The student is correct in principle (assuming cable
is rated for $13\,\text{A}$ continuous). However the
assumption requires that all appliances run simultaneously
and are resistive loads. In practice, motors (AC, fridge)
have higher start-up currents; inductive loads may have
higher peak currents. The $3120\,\text{W}$ is the
\textit{maximum} for a $13\,\text{A}$ rated supply
at $240\,\text{V}$.

% ============================================================
%  CHAPTER 25
% ============================================================
\endgroup

\AnswersChapterBand{25}{Magnetism and Electromagnetism}

\AnswersSectionHeading{Section A}
\begin{multicols}{2}
\begingroup
\let\textbf\AnswersTextbf
\def\quad{ --- }
\setlength{\parindent}{0pt}

\textbf{A1.~B} \quad $B$ at centre of circular coil
$=\mu_0NI/(2r)$.

\textbf{A2.~C} \quad $F=BIL\sin\theta$; wire parallel
to field: $\theta=0°$, $F=0$.

\textbf{A3.~A} \quad $F=Bqv=0.2\times1.6\times10^{-19}
\times10^6=3.2\times10^{-14}\,\text{N}$.

\textbf{A4.~B} \quad Parallel currents (same direction)
attract each other.

\textbf{A5.~B} \quad Commutator reverses current in
coil each half turn, keeping torque in the same
direction.

\textbf{A6.~C} \quad $r=mv/(Bq)$; larger momentum
$\Rightarrow$ larger radius.

\textbf{A7.~B} \quad Hall effect: measures $B$ and
identifies the sign of majority charge carriers.

\textbf{A8.~B} \quad Shunt $S$ in parallel with
galvanometer:
$I_gG=(I-I_g)S\Rightarrow S=GI_g/(I-I_g)$.

\textbf{A9.~A} \quad $B=\mu_0nI=4\pi\times10^{-7}
\times500\times2=4\pi\times10^{-4}
\approx1.26\times10^{-3}\,\text{T}$.

\textbf{A10.~B} \quad Fleming's Left-Hand Rule:
force on current-carrying conductor (motor effect).

\endgroup
\end{multicols}

\AnswersSectionHeading{Section B}
\begingroup
\let\textbf\AnswersTextbf
\setlength{\parindent}{0pt}
\def\quad{\par\smallskip}
\catcode`(=\active
\def({\AnswersParen}

\textbf{B1.}
(a)~Length $=0.08\,\text{m}$:
$F=NBIL\sin90°=50\times0.30\times2.0\times0.08
=2.4\,\text{N}$ (on each of the two longer sides,
same magnitude, opposite direction). \quad
(b)~Torque $=F\times d=2.4\times0.05=0.12\,\text{N\,m}$
(where $d=0.05\,\text{m}$ is half the shorter side
for a flat coil perpendicular to field, i.e.\ 
torque $=NBAI=50\times0.3\times(0.08\times0.05)\times2
=0.12\,\text{N\,m}$). \quad
(c)~$\tau=BINA\cos\alpha=0.30\times2\times50
\times0.004\times\cos30°=0.12\times0.866
=0.104\,\text{N\,m}$. \quad
(d)~Maximum torque when coil is parallel to field
($\alpha=0°$); zero torque when coil is perpendicular
to field ($\alpha=90°$, ``dead point''). At dead point
the coil must have enough momentum to carry it through;
the commutator must switch current at this moment.

\textbf{B2.}
(a)~$KE=qV=2\times1.6\times10^{-19}\times500
=1.6\times10^{-16}\,\text{J}$. \quad
(b)~$v=\sqrt{2KE/m}=\sqrt{2\times1.6\times10^{-16}/
6.64\times10^{-27}}=\sqrt{4.82\times10^{10}}
=2.19\times10^5\,\text{m\,s}^{-1}$. \quad
(c)~$r=mv/(Bq)=6.64\times10^{-27}
\times2.19\times10^5/(0.20\times2\times1.6\times10^{-19})
=1.454\times10^{-21}/6.4\times10^{-20}
=0.0227\,\text{m}=2.27\,\text{cm}$. \quad
(d)~$T=2\pi m/(Bq)=2\pi\times6.64\times10^{-27}/
(0.20\times3.2\times10^{-19})=4.17\times10^{-26}/
6.4\times10^{-20}=6.52\times10^{-7}\,\text{s}$. \quad
(e)~Circular path in $xy$-plane; after one complete
orbit, particle returns to origin $(0,0,0)$.

\textbf{B3.}
(a)~$V_H=BI/(nqw)$:
$B=V_H nqw/I=3.5\times10^{-3}\times2\times10^{21}
\times1.6\times10^{-19}\times2\times10^{-3}/
(50\times10^{-3})
=3.5\times10^{-3}\times6.4\times10^{-1}/0.05
=3.5\times10^{-3}\times12.8=0.0448\,\text{T}
\approx45\,\text{mT}$. \quad
(b)~With current parallel to field: force on carriers
$=qv\times B=0$ (no transverse deflection);
$V_H=0$. \quad
(c)~If carriers are positive: they accumulate on one
face; Hall voltage $+$ on that face. If negative
(electrons): they accumulate on the opposite face;
$+$ on the other face. Sign of $V_H$ reveals carrier
type. \quad
(d)~$V_H=BI/(nqw)=0.0448\times50\times10^{-3}/
(8.5\times10^{28}\times1.6\times10^{-19}
\times2\times10^{-3})
=2.24\times10^{-3}/(2.72\times10^7)
=8.2\times10^{-11}\,\text{V}$. Tiny because $n$ in
copper is $\sim10^7$ times larger than in the
semiconductor; more carriers mean each carries less
of the deflecting force, giving a much smaller
transverse voltage.

\textbf{B4.}
(a)~$F/L=\mu_0I_AI_B/(2\pi d)
=4\pi\times10^{-7}\times6\times4/(2\pi\times0.04)
=9.6\times10^{-6}/0.04\pi/(2\pi)$... let me compute directly:
$F/L=2\times10^{-7}\times6\times4/0.04
=2\times10^{-7}\times600=1.2\times10^{-4}\,\text{N\,m}^{-1}$.
Same direction $\Rightarrow$ \textbf{attractive}. \quad
(b)~Wire C is at $d=0.02\,\text{m}$ from each of A
and B (midway). Force from A (same direction?):
A carries $6\,\text{A}$ (same direction as C? C is
opposite direction to both A and B).
$F_{A\to C}/L=2\times10^{-7}\times6\times3/0.02
=1.8\times10^{-4}\,\text{N\,m}^{-1}$ (repulsive
from A, toward B).
$F_{B\to C}/L=2\times10^{-7}\times4\times3/0.02
=1.2\times10^{-4}\,\text{N\,m}^{-1}$ (repulsive
from B, toward A).
Net force on C $=(1.8-1.2)\times10^{-4}
=6\times10^{-5}\,\text{N\,m}^{-1}$ toward B. \quad
(c)~Zero field between A and B: Let distance from
A be $x$.
$\mu_0I_A/(2\pi x)=\mu_0I_B/(2\pi(0.04-x))$;
$6/x=4/(0.04-x)$; $6(0.04-x)=4x$;
$0.24=10x$; $x=0.024\,\text{m}=2.4\,\text{cm}$ from A.

\textbf{B5.}
(a)~\textit{(Ray diagram description)}:
Rectangular coil between N and S pole pieces;
commutator (split ring) contacts two brushes;
current enters through one brush, flows through coil,
exits through other brush; forces on sides perpendicular
to field are upward and downward (by FLHR),
creating torque. \quad
(b)~$\tau_{max}=BINA=0.15\times3\times200
\times(20\times10^{-4})=0.18\,\text{N\,m}$.
\quad
(c)~(i)~$\mathcal{E}_{back}=V-IR_{coil}
=12-2\times1.5=12-3=9\,\text{V}$.
(ii)~At start: back-EMF $=0$; $I_{start}=V/R
=12/1.5=8\,\text{A}$, which may exceed the rated
current and burn out the armature winding.
(iii)~$P_{input}=VI=12\times2=24\,\text{W}$;
$P_{coil}=I^2R=4\times1.5=6\,\text{W}$;
$P_{mech}=P_{input}-P_{coil}=24-6=18\,\text{W}$.
(iv)~$\eta=P_{mech}/P_{input}=18/24=75\%$.

% ============================================================
%  CHAPTER 26
% ============================================================
\endgroup

\AnswersChapterBand{26}{Electromagnetic Induction}

\AnswersSectionHeading{Section A}
\begin{multicols}{2}
\begingroup
\let\textbf\AnswersTextbf
\def\quad{ --- }
\setlength{\parindent}{0pt}

\textbf{A1.~B} \quad The induced EMF equals the rate of
change of flux linkage (Faraday's Law); it is greatest
when $d\Phi/dt$ is maximum, not when $\Phi$ itself is
large.

\textbf{A2.~C} \quad By Lenz's Law, eddy currents induced
in the copper tube create a magnetic braking force
opposing the magnet's motion. The magnet falls slower
than free fall --- the classic ``falling magnet''
demonstration of electromagnetic braking.

\textbf{A3.~C} \quad $\mathcal{E}=BLv
=0.25\times0.40\times15=1.5\,\text{V}$.

\textbf{A4.~C} \quad Lenz's Law states the induced
current opposes the change causing it. If it aided the
change, it would amplify itself indefinitely ---
creating energy from nothing. Lenz's Law is therefore
a direct consequence of conservation of energy.

\textbf{A5.~B} \quad Peak EMF $\mathcal{E}_0=NBA\omega$.
Doubling $N$ while keeping $B$, $A$ and $\omega$
constant doubles $\mathcal{E}_0$. Option A gives
$\tfrac{1}{2}N\times2\omega=N\omega$ (unchanged).

\textbf{A6.~C} \quad Laminations divide the iron core
into thin, electrically insulated sheets. Eddy currents
are confined to narrow strips of high resistance,
reducing $I^2R$ losses dramatically.

\textbf{A7.~C} \quad Energy stored in an inductor is
$W=\tfrac{1}{2}LI^2$, analogous to kinetic energy
$\tfrac{1}{2}mv^2$.

\textbf{A8.~C} \quad Back-EMF $\propto$ angular speed.
At full speed the motor rotates fastest;
back-EMF is maximum, limiting the current to its
lowest running value.

\textbf{A9.~C} \quad $|\mathcal{E}|=N\Delta\Phi/\Delta t
=100\times(0.08-0.02)/0.05
=100\times1.2=120\,\text{V}$.

\textbf{A10.~B} \quad Slip rings maintain continuous
sliding contact as the coil rotates, delivering the
alternating current without reversing it. A commutator
would reverse the current every half-turn, producing
pulsating DC.

\endgroup
\end{multicols}

\AnswersSectionHeading{Section B}
\begingroup
\let\textbf\AnswersTextbf
\setlength{\parindent}{0pt}
\def\quad{\par\smallskip}
\catcode`(=\active
\def({\AnswersParen}

\textbf{B1.}
(a)~$A=\pi r^2=\pi(0.04)^2=5.027\times10^{-3}\,\text{m}^2$.
$\Phi_i=B_iA=0.60\times5.027\times10^{-3}
=3.016\times10^{-3}\,\text{Wb}$.
$\Phi_f=0.10\times5.027\times10^{-3}
=5.027\times10^{-4}\,\text{Wb}$. \quad
(b)~$|\mathcal{E}|=N\Delta\Phi/\Delta t
=150\times(3.016-0.503)\times10^{-3}/0.20
=150\times2.513\times10^{-3}/0.20
=1.88\,\text{V}$. \quad
(c)~$I=\mathcal{E}/R=1.88/5.0=0.377\,\text{A}
\approx0.38\,\text{A}$. \quad
(d)~Flux is decreasing. By Lenz's Law the induced
current must \textit{oppose} this decrease --- it creates
a magnetic field in the \textbf{same direction} as the
original (decreasing) field, trying to maintain the
flux. Using the right-hand rule, the current circulates
in the same rotational sense as the current that
originally produced the field direction. \quad
(e)~Coil rotated $90°$: flux changes from
$\Phi_i=3.016\times10^{-3}\,\text{Wb}$ (coil
perpendicular to $B$) to $\Phi_f=0$ (coil parallel to
$B$).
$|\mathcal{E}|=N\Delta\Phi/\Delta t
=150\times3.016\times10^{-3}/0.20=2.26\,\text{V}$.

\textbf{B2.}
(a)~$\mathcal{E}=BLv=0.80\times0.50\times2.5
=1.0\,\text{V}$;
$I=\mathcal{E}/R=1.0/0.40=2.5\,\text{A}$. \quad
(b)~Braking force $F=BIL
=0.80\times2.5\times0.50=1.0\,\text{N}$.
At constant velocity the applied force equals the
braking force: $F_{applied}=1.0\,\text{N}$. \quad
(c)~$P_{input}=Fv=1.0\times2.5=2.5\,\text{W}$;
$P_{dissipated}=I^2R=6.25\times0.40=2.5\,\text{W}$.
They are equal: all mechanical power input is
converted to electrical heat in the rod ---
energy is conserved. The rod is a linear generator
operating at $100\%$ conversion into electrical
power (all of which is dissipated in $R$). \quad
(d)~Equation of motion when external force removed:
$ma=-B^2L^2v/R$ (braking force $\propto v$).
As $v$ decreases, the braking force decreases
proportionally; deceleration therefore decreases.
The rod asymptotically approaches zero velocity
--- it never stops abruptly but decelerates
ever more slowly. \quad
(e)~$v=\mathcal{E}/(BL)=1.0/(0.80\times0.50)
=2.5\,\text{m\,s}^{-1}$
(confirming part (a)).

\textbf{B3.}
(a)~$\omega=2\pi f=2\pi\times25=157\,\text{rad\,s}^{-1}$.
\quad
(b)~$\mathcal{E}_0=NBA\omega
=500\times0.20\times80\times10^{-4}\times157
=500\times0.20\times0.0080\times157=125.6\,\text{V}
\approx126\,\text{V}$. \quad
(c)~$\mathcal{E}_{rms}=\mathcal{E}_0/\sqrt{2}
=126/1.414=89.1\,\text{V}\approx89\,\text{V}$. \quad
(d)~Total resistance $=R_{coil}+R_L=2+50=52\,\Omega$.
Peak current $I_0=\mathcal{E}_0/R_{total}
=126/52=2.42\,\text{A}$.
RMS current $=2.42/\sqrt{2}=1.71\,\text{A}$.
Power to load $=I_{rms}^2\times R_L
=2.924\times50=146\,\text{W}$. \quad
(e)~By Lenz's Law, the induced current in the rotating
coil experiences a force opposing rotation ---
an electromagnetic braking torque. To maintain constant
angular speed, an external mechanical torque must
continuously supply power equal to the electrical power
generated. This is how mechanical energy (from turbines,
etc.) is converted to electrical energy. \quad
(f)~Doubling frequency doubles $\omega$; since
$\mathcal{E}_0\propto\omega$, the new peak EMF
$=2\times126=252\,\text{V}$.

\textbf{B4.}
(a)~$L=\mu_0N^2A/\ell
=4\pi\times10^{-7}\times(600)^2
\times8.0\times10^{-4}/0.25$.
$=1.2566\times10^{-6}\times3.6\times10^5
\times3.2\times10^{-3}
=1.447\times10^{-3}\,\text{H}\approx1.45\,\text{mH}$.
\quad
(b)~$\mathcal{E}=-L\,dI/dt
=1.447\times10^{-3}\times4.0/0.050
=1.447\times10^{-3}\times80=0.116\,\text{V}$. \quad
(c)~$W=\tfrac{1}{2}LI^2
=\tfrac{1}{2}\times1.447\times10^{-3}\times16
=11.6\times10^{-3}\,\text{J}=11.6\,\text{mJ}$. \quad
(d)~$\mathcal{E}_2=M\,dI_1/dt
=0.72\times10^{-6}\times200
=1.44\times10^{-4}\,\text{V}=0.144\,\text{mV}$. \quad
(e)~An iron core has relative permeability
$\mu_r\gg1$ (typically hundreds to thousands).
Since $L=\mu_r\mu_0N^2A/\ell$, the inductance
scales directly with $\mu_r$. The iron core channels
far more magnetic flux through the solenoid per unit
current, increasing the flux linkage $N\Phi$ and hence
$L$ by the same factor $\mu_r$.

\textbf{B5.}
(a)~$\mathcal{E}_{back}=V-IR_{coil}
=240-3.0\times2.0=234\,\text{V}$. \quad
(b)~$P_{input}=VI=240\times3.0=720\,\text{W}$;
$P_{heat}=I^2R=(3.0)^2\times2.0=18\,\text{W}$;
$P_{mechanical}=720-18=702\,\text{W}$. \quad
(c)~$\eta=P_{mech}/P_{input}=702/720=97.5\%$. \quad
(d)~At startup, $\mathcal{E}_{back}=0$ (coil not
spinning); $I_{start}=V/R_{coil}=240/2.0=120\,\text{A}$.
This is $40\times$ the normal running current.
Such a large current can burn out armature windings,
melt commutator brushes, or blow circuit protection,
so direct-on-line starting of large motors is dangerous
without a starter resistor. \quad
(e)~With $18\,\Omega$ in series:
$I_{start}=240/(2.0+18)=240/20=12\,\text{A}$.
$P_{starter}=I^2\times18=144\times18=2592\,\text{W}$.
Remove starter resistor when back-EMF is large enough
to keep $I\leq6.0\,\text{A}$:
$240-\mathcal{E}_{back}=6.0\times2.0=12\,\text{V}$;
$\mathcal{E}_{back}=228\,\text{V}$.
Since $\mathcal{E}_{back}\propto$ speed:
fraction of full speed $=228/234\approx97.4\%$.
Remove the starter resistor at approximately
$97\%$ of full speed.

% ============================================================
%  CHAPTER 27
% ============================================================
\endgroup

\AnswersChapterBand{27}{Atomic Structure and Spectra}

\AnswersSectionHeading{Section A}
\begin{multicols}{2}
\begingroup
\let\textbf\AnswersTextbf
\def\quad{ --- }
\setlength{\parindent}{0pt}

\textbf{A1.~C} \quad Rutherford's back-scattered alpha
particles implied all the positive charge and almost all
the mass are concentrated in an incredibly small region
--- the nucleus. The atom is predominantly empty space
with electrons orbiting far outside.

\textbf{A2.~C} \quad $E_4=-13.6/4^2=-13.6/16
=-0.85\,\text{eV}$.

\textbf{A3.~B} \quad The minimum frequency required to
eject photoelectrons is the \textbf{threshold frequency}
$f_0=\phi/h$.

\textbf{A4.~C} \quad Each photon interacts with one
electron. Greater intensity means more photons per
second $\Rightarrow$ more electrons ejected per second.
The energy per photon (and hence $KE_{max}$) depends
only on frequency, which is unchanged.

\textbf{A5.~B} \quad De Broglie's relation:
$\lambda=h/p$.

\textbf{A6.~B} \quad $\Delta E=E_4-E_2
=-0.85-(-3.40)=+2.55\,\text{eV}$.

\textbf{A7.~C} \quad The stopping potential $V_s$ is
defined by $eV_s=KE_{max}$, so $V_s=KE_{max}/e$.
It directly measures the maximum kinetic energy of
the emitted electrons.

\textbf{A8.~C} \quad The photoelectric effect --- in
particular the existence of a threshold frequency and
the immediate emission even in faint light --- cannot
be explained by wave theory. Einstein's photon
explanation (quantised light packets) provides
direct evidence for the particle nature of light.

\textbf{A9.~A} \quad The Lyman series consists of all
transitions ending at $n=1$ (the ground state).

\textbf{A10.~B} \quad The Heisenberg Uncertainty
Principle is a fundamental property of quantum
systems, not a technological limitation. A particle
does not simultaneously possess both a definite
position and a definite momentum; the more precisely
one is known, the less precisely the other can be,
regardless of how good the measuring instrument is.

\endgroup
\end{multicols}

\AnswersSectionHeading{Section B}
\begingroup
\let\textbf\AnswersTextbf
\setlength{\parindent}{0pt}
\def\quad{\par\smallskip}
\catcode`(=\active
\def({\AnswersParen}

\textbf{B1.}
(a)~$7.7\,\text{MeV}=7.7\times10^6\times1.6
\times10^{-19}=1.232\times10^{-12}\,\text{J}$. \quad
(b)~$q_\alpha=2e=3.2\times10^{-19}\,\text{C}$;
$q_{Au}=79e=1.264\times10^{-17}\,\text{C}$.
$d=kZq_\alpha q_{Au}/KE
=(9\times10^9\times3.2\times10^{-19}
\times1.264\times10^{-17})/1.232\times10^{-12}
=(3.640\times10^{-26})/1.232\times10^{-12}
=2.95\times10^{-14}\,\text{m}$.
The nuclear radius is $\sim7\times10^{-15}\,\text{m}$;
the distance of closest approach
($\approx4$ nuclear radii) confirms the alpha particle
just barely probes the nuclear surface in a
head-on collision. \quad
(c)~Atomic radius $\sim10^{-10}\,\text{m}$;
nuclear radius $\sim10^{-15}\,\text{m}$.
The ratio of volumes is $(10^{-10})^3/(10^{-15})^3
=10^{15}$. The nucleus occupies roughly $10^{-15}$
of the atomic volume. An alpha particle aimed at a
random point has a $\sim1$ in $10^{15}$ chance of
intercepting the nucleus, which is why most pass
straight through. \quad
(d)~The nucleus occupies a volume fraction of about
$10^{-15}$ of the total atomic volume.
The remaining $\sim99.9999999999999\%$ contains
only the diffuse electron cloud, which is too light
to deflect a fast alpha particle significantly.
The atom is overwhelmingly empty space.

\textbf{B2.}
(a)~$E=hc/\lambda
=(6.63\times10^{-34}\times3.0\times10^8)/97.2
\times10^{-9}=2.046\times10^{-18}\,\text{J}
=2.046\times10^{-18}/1.6\times10^{-19}
=12.8\,\text{eV}$. \quad
(b)~$\Delta E=E_n-E_1=12.8\,\text{eV}$;
$E_1=-13.6\,\text{eV}$; $E_n=12.8-13.6
=-0.80\,\text{eV}\approx E_4=-0.85\,\text{eV}$.
The transition is $n=1\to n=4$ (Lyman series
absorption). \quad
(c)~$n=4\to n=2$:
$\Delta E=-0.85-(-3.40)=2.55\,\text{eV}
=4.08\times10^{-19}\,\text{J}$.
$\lambda=hc/\Delta E=1.989\times10^{-25}/4.08
\times10^{-19}=4.88\times10^{-7}\,\text{m}
=488\,\text{nm}$. \quad
(d)~$488\,\text{nm}$ is within the visible range
(400--700\,nm). It appears \textbf{blue-green (cyan)}
--- this is the H$\beta$ Balmer line. \quad
(e)~Ionisation from ground state: $E=13.6\,\text{eV}
=2.176\times10^{-18}\,\text{J}$.
$\lambda=hc/E=1.989\times10^{-25}/2.176
\times10^{-18}=9.14\times10^{-8}\,\text{m}
=91.4\,\text{nm}$ (UV; the Lyman limit).

\textbf{B3.}
(a)~$f_0=\phi/h
=(4.3\times1.6\times10^{-19})/6.63\times10^{-34}
=6.88\times10^{-19}/6.63\times10^{-34}
=1.038\times10^{15}\,\text{Hz}$.
$\lambda_0=c/f_0=3.0\times10^8/1.038\times10^{15}
=2.89\times10^{-7}\,\text{m}=289\,\text{nm}$
(UV). \quad
(b)~$E_{photon}=hc/\lambda=1.989\times10^{-25}/200
\times10^{-9}=9.945\times10^{-19}\,\text{J}
=6.22\,\text{eV}$.
$KE_{max}=6.22-4.3=1.92\,\text{eV}
=3.072\times10^{-19}\,\text{J}$.
$v_{max}=\sqrt{2KE/m_e}
=\sqrt{2\times3.072\times10^{-19}/9.11\times10^{-31}}
=\sqrt{6.745\times10^{11}}
=8.21\times10^5\,\text{m\,s}^{-1}$. \quad
(c)~$V_s=KE_{max}/e=1.92\,\text{V}$. \quad
(d)~(i)~$KE_{max}$ is \textbf{unchanged}: it depends
only on photon frequency, which is unchanged by
increasing intensity.
(ii)~The rate of electron emission \textbf{doubles}:
twice as many photons per second means twice as many
photoelectron ejections per second. \quad
(e)~$f_0=c/\lambda_0=3.0\times10^8/500\times10^{-9}
=6.0\times10^{14}\,\text{Hz}$.
$\phi=hf_0=6.63\times10^{-34}\times6.0\times10^{14}
=3.978\times10^{-19}\,\text{J}=2.49\,\text{eV}$.
Zinc ($\phi=4.3\,\text{eV}$) requires higher-energy
(shorter-wavelength) photons. This second metal has a
lower work function and is easier to ionise
photelectrically.

\textbf{B4.}
(a)~$KE=eV=\tfrac{1}{2}m_ev^2\Rightarrow
v=\sqrt{2eV/m_e}$; $p=m_ev=\sqrt{2m_eeV}$;
$\lambda=h/p=h/\sqrt{2m_eeV}$. \checkmark \quad
(b)~$V=50\,\text{V}$:
$\lambda=6.63\times10^{-34}/\sqrt{2\times9.11
\times10^{-31}\times1.6\times10^{-19}\times50}
=6.63\times10^{-34}/3.82\times10^{-24}
=1.74\times10^{-10}\,\text{m}=0.174\,\text{nm}$.
$V=200\,\text{V}$: $\lambda\propto1/\sqrt{V}$;
$\lambda=0.174\times\sqrt{50/200}=0.087\,\text{nm}$.
$V=10\,000\,\text{V}$: $\lambda=0.174
\times\sqrt{50/10000}=0.174\times0.0707
=0.0123\,\text{nm}$. \quad
(c)~Crystal lattice spacings $0.1$--$0.5\,\text{nm}$.
At $50\,\text{V}$ ($\lambda=0.174\,\text{nm}$) and
$200\,\text{V}$ ($\lambda=0.087\,\text{nm}$),
$\lambda$ is comparable to lattice spacings
$\Rightarrow$ diffraction expected.
At $10\,000\,\text{V}$ ($\lambda=0.012\,\text{nm}$),
$\lambda$ is much smaller than lattice spacings
$\Rightarrow$ diffraction is negligible. \quad
(d)~Visible light: 400--700\,nm. Electron wavelengths
at these voltages are $10^3$--$10^4$ times smaller.
Since resolution limit $\approx\lambda/2$, electron
microscopes resolve features thousands of times smaller
than optical microscopes --- allowing imaging of
individual atoms. \quad
(e)~$\lambda_p=h/\sqrt{2m_peV}
=6.63\times10^{-34}/\sqrt{2\times1.67\times10^{-27}
\times1.6\times10^{-19}\times200}
=6.63\times10^{-34}/3.27\times10^{-22}
=2.03\times10^{-12}\,\text{m}=0.00203\,\text{nm}$.
Since $\lambda\propto1/\sqrt{m}$ (same $V$):
$\lambda_p/\lambda_e=\sqrt{m_e/m_p}
=\sqrt{9.11\times10^{-31}/1.67\times10^{-27}}
=1/43$.
The proton's de Broglie wavelength is 43 times
smaller than the electron's at the same accelerating
voltage, because the proton acquires the same kinetic
energy but has far greater mass and therefore far
greater momentum.

\textbf{B5.}
(a)~Centripetal force $=$ Coulomb force:
$m_ev^2/r=ke^2/r^2\Rightarrow v=\sqrt{ke^2/(m_er)}$.
$v=\sqrt{(9\times10^9\times(1.6\times10^{-19})^2)/
(9.11\times10^{-31}\times5.29\times10^{-11})}
=\sqrt{2.304\times10^{-28}/4.819\times10^{-41}}
=\sqrt{4.782\times10^{12}}
=2.19\times10^6\,\text{m\,s}^{-1}$. \quad
(b)~$p=m_ev=9.11\times10^{-31}\times2.19\times10^6
=1.995\times10^{-24}\,\text{kg\,m\,s}^{-1}$.
$\lambda=h/p=6.63\times10^{-34}/1.995\times10^{-24}
=3.32\times10^{-10}\,\text{m}$.
Circumference $=2\pi r_1
=2\pi\times5.29\times10^{-11}
=3.32\times10^{-10}\,\text{m}$.
$\lambda=2\pi r_1$ exactly: one complete de Broglie
wavelength fits around the Bohr orbit. This is the
standing-wave condition that gives physical meaning to
Bohr's quantisation postulate. \quad
(c)~Classical radiated power:
$P=e^2a^2/(6\pi\varepsilon_0c^3)$.
$=(1.6\times10^{-19})^2\times(9.0\times10^{22})^2/
(6\pi\times8.85\times10^{-12}\times(3\times10^8)^3)$.
Numerator: $2.56\times10^{-38}\times8.1\times10^{45}
=2.074\times10^8$.
Denominator: $6\pi\times8.85\times10^{-12}
\times2.7\times10^{25}=4.506\times10^{15}$.
$P=2.074\times10^8/4.506\times10^{15}
\approx4.6\times10^{-8}\,\text{W}$.
The electron's kinetic energy $=\tfrac{1}{2}m_ev^2
\approx2.18\times10^{-18}\,\text{J}$.
Time to radiate away all kinetic energy
$\approx2.18\times10^{-18}/4.6\times10^{-8}
\approx4.7\times10^{-11}\,\text{s}$.
Classical physics predicts atomic collapse in
$\sim10^{-11}\,\text{s}$. Atoms are stable because
quantum mechanics forbids the electron from losing
energy below the ground state; the stationary states
are not arbitrary postulates but consequences of the
wave nature of the electron. \quad
(d)~Bohr arbitrarily \textit{postulated} quantised
orbits without physical justification --- it was known
even then that an accelerating classical charge must
radiate. The model works only for hydrogen and
hydrogen-like ions, cannot predict spectral line
intensities, fails completely for multi-electron atoms,
and violates the uncertainty principle (precise orbits
imply simultaneously exact position and momentum).
It was a phenomenological model that produced correct
energy levels for the wrong physical reasons, and its
incompleteness was recognised from the outset.

% ============================================================
%  CHAPTER 28
% ============================================================
\endgroup

\AnswersChapterBand{28}{Radioactivity and Nuclear Physics}

\AnswersSectionHeading{Section A}
\begin{multicols}{2}
\begingroup
\let\textbf\AnswersTextbf
\def\quad{ --- }
\setlength{\parindent}{0pt}

\textbf{A1.~B} \quad An alpha particle has the same
composition as a helium-4 nucleus: 2 protons and
2 neutrons, ${}^4_2\text{He}$.

\textbf{A2.~C} \quad In $\beta^-$ decay a neutron in
the nucleus converts to a proton, emitting a fast
electron (beta particle) and an antineutrino. The
proton number $Z$ increases by 1; $A$ is unchanged.

\textbf{A3.~C} \quad $1\,\text{h}=60\,\text{min}
=3\times t_{1/2}$. Fraction remaining
$=(1/2)^3=1/8$.

\textbf{A4.~C} \quad Gamma radiation (high-energy
photons) is the most penetrating: it requires several
centimetres of lead or metres of concrete to attenuate
significantly and is never completely stopped.

\textbf{A5.~C} \quad Binding energy per nucleon peaks
at ${}^{56}_{26}$Fe ($\approx8.8\,\text{MeV}$ per
nucleon), making iron-56 the most stable nucleus.

\textbf{A6.~B} \quad The moderator slows fast
(``fast'') neutrons released in fission to thermal
speeds. Slow (thermal) neutrons are much more
effective at inducing further fission in ${}^{235}$U.

\textbf{A7.~C} \quad The fused product lies higher on
the binding-energy-per-nucleon curve than the light
reactants. Greater binding energy per nucleon means
less mass per nucleon; the mass deficit is released as
energy via $E=mc^2$.

\textbf{A8.~C} \quad Alpha decay: $Z\to Z-2$.
Each $\beta^-$ decay: $Z\to Z+1$.
Net after $\alpha$ then $2\beta^-$:
$Z-2+2=Z$. Final proton number unchanged.

\textbf{A9.~C} \quad $\lambda=\ln2/t_{1/2}
=0.693/t_{1/2}$.

\textbf{A10.~B} \quad ${}^{14}$C has a half-life of
5730 years, which is comparable to the timescale of
human history and prehistoric archaeology (hundreds
to tens of thousands of years). A too-short half-life
would leave nothing detectable; a too-long one would
show negligible change.

\endgroup
\end{multicols}

\AnswersSectionHeading{Section B}
\begingroup
\let\textbf\AnswersTextbf
\setlength{\parindent}{0pt}
\def\quad{\par\smallskip}
\catcode`(=\active
\def({\AnswersParen}

\textbf{B1.}
(a)~$\Delta A=238-206=32$; each $\alpha$ reduces
$A$ by 4; number of $\alpha$ decays $=32/4=8$.
Net $Z$ change from 8 $\alpha$ decays:
$92-16=76$; but $Z_{Pb}=82$; deficit $=82-76=6$;
each $\beta^-$ increases $Z$ by 1; so 6 $\beta^-$
decays. \textbf{8 alpha and 6 beta-minus decays.}
Check: $A=238-32=206$ \checkmark;
$Z=92-16+6=82$ \checkmark. \quad
(b)~${}^{238}_{92}\text{U}\to
{}^{234}_{90}\text{Th}+{}^{4}_{2}\text{He}$. \quad
(c)~First $\beta^-$:
${}^{234}_{90}\text{Th}\to
{}^{234}_{91}\text{Pa}+{}^{0}_{-1}e+\bar{\nu}_e$.
Second $\beta^-$:
${}^{234}_{91}\text{Pa}\to
{}^{234}_{92}\text{U}+{}^{0}_{-1}e+\bar{\nu}_e$. \quad
(d)~$t_{1/2}=4.47\times10^9\times365.25\times86400
=1.411\times10^{17}\,\text{s}$.
$\lambda=0.693/1.411\times10^{17}
=4.91\times10^{-18}\,\text{s}^{-1}$. \quad
(e)~$N=(1.00/238)\times6.02\times10^{23}
=2.529\times10^{21}\,\text{atoms}$.
$A=\lambda N=4.91\times10^{-18}\times2.529
\times10^{21}=1.24\times10^4\,\text{Bq}$.
($1.00\,\text{g}$ of ${}^{238}$U has an activity of
only about 12\,400\,Bq --- U-238 decays very slowly.)

\textbf{B2.}
(a)~$t_{1/2}=8.04\times86400=6.947\times10^5\,\text{s}$.
$\lambda=0.693/6.947\times10^5
=9.976\times10^{-7}\,\text{s}^{-1}
\approx9.98\times10^{-7}\,\text{s}^{-1}$. \quad
(b)~$N_0=A_0/\lambda=3.7\times10^9/9.976
\times10^{-7}=3.71\times10^{15}\,\text{atoms}$. \quad
(c)~$m=N_0\times M/N_A
=3.71\times10^{15}\times131/6.02\times10^{23}
=8.07\times10^{-7}\,\text{g}=0.807\,\mu\text{g}$. \quad
(d)~$32.16\,\text{days}=32.16/8.04
=4.00\times t_{1/2}$.
$A=3.7\times10^9\times(1/2)^4
=3.7\times10^9/16=2.31\times10^8\,\text{Bq}$. \quad
(e)~Decay curve: exponential starting at
$3.7\times10^9\,\text{Bq}$, halving every 8.04 days.
Label axes ``Activity (Bq)'' and ``Time (days)''.
Mark: $t=0$: $3.70\times10^9$;
$t=8.04$: $1.85\times10^9$;
$t=16.08$: $9.25\times10^8$;
$t=24.12$: $4.63\times10^8\,\text{Bq}$. \quad
(f)~Activity decreases exponentially with time.
Delaying administration means the sample reaches the
patient with lower activity than prescribed. Since
the therapeutic effect depends on the total radiation
dose delivered to the thyroid, a lower-activity sample
may fail to achieve the required dose. The dose is
prescribed and calibrated for a specific activity at
the time of administration, not at the time of
preparation.

\textbf{B3.}
(a)~${}^{56}_{26}$Fe has 26 protons and 30 neutrons.
Mass of constituents:
$26\times1.007276+30\times1.008665
=26.18918+30.25995=56.44913\,\text{u}$.
Mass defect:
$\Delta m=56.44913-55.934939=0.51419\,\text{u}$.
Binding energy:
$E_B=0.51419\times931.5=479.0\,\text{MeV}$. \quad
(b)~$E_B/A=479.0/56=8.55\,\text{MeV\,per\,nucleon}$. \quad
(c)~${}^4$He: $E_B/A=28.3/4=7.08\,\text{MeV\,per\,nucleon}$.
${}^{56}$Fe has higher binding energy per nucleon
($8.55$ vs $7.08\,\text{MeV}$) --- it is more tightly
bound and therefore more stable. \quad
(d)~(i)~${}^{235}$U lies to the right of the peak
(lower $E_B/A$). Fission fragments (mass number
$\sim100$--$140$) sit nearer the peak where $E_B/A$
is higher. Products are more tightly bound $\Rightarrow$
less total mass $\Rightarrow$ mass deficit released as
energy.
(ii)~Hydrogen isotopes lie far to the left of the peak
(low $E_B/A$). Fusion to helium or heavier nuclei
moves up the curve to higher $E_B/A$ $\Rightarrow$
mass deficit released as energy.
(iii)~${}^{56}$Fe is \textit{at} the peak. Any
reaction (fission or fusion) moves to a lower point on
the curve $\Rightarrow$ products have lower $E_B/A$,
meaning more mass must be supplied rather than
released. No energy is liberated.

\textbf{B4.}
(a)~Fraction remaining
$=A/A_0=1.5/6.5=0.2308$. \quad
(b)~$t=\frac{t_{1/2}}{\ln2}\ln\!\left(\frac{A_0}{A}\right)
=\frac{5730}{0.693}\ln\!\left(\frac{6.5}{1.5}\right)
=8268\times\ln(4.333)
=8268\times1.466=12\,121\,\text{years}
\approx12\,100\,\text{years}$. \quad
(c)~Upper bound ($A=1.3$): $t=8268\times\ln(6.5/1.3)
=8268\times1.609=13\,305\,\text{years}$.
Lower bound ($A=1.7$): $t=8268\times\ln(6.5/1.7)
=8268\times1.341=11\,087\,\text{years}$.
Uncertainty $\approx\pm1100\,\text{years}$. \quad
(d)~After 50\,000 years: $t/t_{1/2}=50000/5730=8.73$
half-lives. Remaining fraction
$=(0.5)^{8.73}\approx0.0023$ (0.23\% of original).
The activity is barely above background radiation,
making accurate measurement impossible and the
calculated age unreliable. \quad
(e)~$12100$ years places the piece at approximately
$10\,100$ BCE --- well within the Later Stone Age
of West Africa. If the age were 1200 years (the
example given), it would be approximately $826\,\text{CE}$,
coinciding with the early flourishing of Yoruba
kingdoms and the Ife civilisation, which was
producing sophisticated terracotta and later bronze
artwork during this period. The Oyo Empire was also
consolidating at this time.

\textbf{B5.}
(a)~Thermal power $=P_{elec}/\eta
=1000\times10^6/0.35=2.857\times10^9\,\text{W}
\approx2.86\,\text{GW}$. \quad
(b)~Energy per fission $=200\times1.6\times10^{-13}
=3.2\times10^{-11}\,\text{J}$.
Fissions per second $=2.857\times10^9/3.2\times10^{-11}
=8.93\times10^{19}\,\text{s}^{-1}$. \quad
(c)~Mass per ${}^{235}$U atom $=235/6.02\times10^{23}
=3.904\times10^{-22}\,\text{g}$.
Mass per second $=8.93\times10^{19}\times3.904
\times10^{-22}=3.49\times10^{-2}\,\text{g\,s}^{-1}$.
Mass per day $=3.49\times10^{-2}\times86400
\approx3010\,\text{g\,day}^{-1}\approx3.01\,\text{kg\,day}^{-1}$.
\quad
(d)~Nuclear: $3.01\times365\approx1099\,\text{kg}
\approx1.1\,\text{tonnes\,year}^{-1}$.
Coal: thermal power $=1000/0.40=2500\,\text{MW}$.
Energy per year $=2.5\times10^9\times3.156
\times10^7=7.89\times10^{16}\,\text{J}$.
Mass of coal $=7.89\times10^{16}/(24\times10^6)
=3.29\times10^9\,\text{kg}=3.29$ million tonnes.
Nuclear fuel is approximately $\mathbf{3}$ \textbf{million
times} more energy-dense per unit mass than coal. \quad
(e)~$0.001=(1/2)^n\Rightarrow n=\ln(1000)/\ln2
=6.908/0.693=9.97\approx10\,\text{half-lives}$.
Time $=10\times29=290\,\text{years}$.
${}^{90}$Sr must be safely isolated for nearly three
centuries before its activity falls to $0.1\%$ of
its initial level --- illustrating the long-term
waste management challenge of nuclear fission power.

% ============================================================
%  CHAPTER 29
% ============================================================
 \endgroup

\AnswersChapterBand{29}{Electronics}

\AnswersSectionHeading{Section A}
\begin{multicols}{2}
\begingroup
\let\textbf\AnswersTextbf
\def\quad{ --- }
\setlength{\parindent}{0pt}

\textbf{A1.~B} \quad $n$-type silicon is formed by
doping with a pentavalent (Group V) element such as
phosphorus or arsenic. The fifth valence electron is
not needed for bonding and is easily freed as a
conduction electron.

\textbf{A2.~C} \quad A forward-biased silicon diode
requires approximately $0.6$--$0.7\,\text{V}$ to
overcome the built-in potential barrier before
significant current flows ($\approx0.3\,\text{V}$
for germanium).

\textbf{A3.~C} \quad $I_C=h_{FE}\times I_B
=200\times25\times10^{-6}=5.0\times10^{-3}\,\text{A}
=5.0\,\text{mA}$.

\textbf{A4.~D} \quad A full-wave bridge rectifier uses
four diodes arranged in a diamond. Two conduct during
the positive half-cycle and the other two during the
negative half-cycle, always directing current through
the load in the same direction.

\textbf{A5.~B} \quad $A_v=-R_f/R_{in}
=-180/20=-9$.

\textbf{A6.~C} \quad In the common-emitter
configuration the output is $180°$ out of phase
with the input (inverted). The voltage gain is
$A_v\approx-R_C/R_E$; the negative sign indicates
inversion.

\textbf{A7.~B} \quad NAND truth: output is LOW only
when \textit{both} inputs are HIGH. With $A=1$ and
$B=0$, the AND result is $0$; NAND gives the inverse:
output $=1$.

\textbf{A8.~B} \quad The depletion region is depleted
of mobile charge carriers (both electrons and holes
have recombined there). It contains fixed
positive ions on the $n$-side and fixed negative
ions on the $p$-side.

\textbf{A9.~C} \quad A photodiode is reverse biased.
Incident photons create electron-hole pairs in the
depletion region; these are swept across the junction
by the built-in electric field, producing a
photocurrent proportional to light intensity.

\textbf{A10.~C} \quad A NOR gate: output is HIGH
only when \textit{all} inputs are LOW (the inverse
of OR, which is LOW only when all inputs are LOW).

\endgroup
\end{multicols}

\AnswersSectionHeading{Section B}
\begingroup
\let\textbf\AnswersTextbf
\setlength{\parindent}{0pt}
\def\quad{\par\smallskip}
\catcode`(=\active
\def({\AnswersParen}

\textbf{B1.}
(a)~Silicon has a band gap of $\approx1.1\,\text{eV}$
separating a full valence band from an empty conduction
band at absolute zero. This is too large for electrons
to cross thermally at low temperatures (hence insulating
at 0\,K) but small enough that at room temperature
a tiny fraction of electrons gain sufficient thermal
energy to bridge the gap, creating electron-hole pairs.
The result is low but measurable conductivity --- neither
as conductive as a metal (no band gap, conduction band
always partially filled) nor as resistive as an insulator
(band gap $>5\,\text{eV}$, essentially no thermal
excitation). \quad
(b)~Intrinsic carrier density: $1.5\times10^{16}\,\text{m}^{-3}$.
Phosphorus donor density: $10^{22}\,\text{m}^{-3}$.
Ratio: $10^{22}/1.5\times10^{16}\approx6.7\times10^5$.
The doping concentration is about 670\,000 times
larger than the intrinsic carrier density; thermally
generated carriers are completely negligible, and the
electrical properties are determined entirely by the
donor electrons. \quad
(c)~Conductivity ratio $=R_{undoped}/R_{doped}
=2300/0.023=10^5$.
Conductivity increased by a factor of
$\mathbf{100\,000}$. \quad
(d)~Each phosphorus donor that loses its extra
electron to the conduction band becomes a fixed
positive ion locked in the crystal lattice. The
number of free electrons equals the number of
fixed positive ions (plus negligible intrinsic
carriers). Total charge
$=$ (mobile electrons) $+$ (fixed $+$ ions) $+$
(intrinsic holes and electrons) $= 0$.
The semiconductor is electrically neutral overall.

\textbf{B2.}
(a)~$V_{peak}=V_{rms}\times\sqrt{2}
=12\times1.414=17.0\,\text{V}$. \quad
(b)~Two diodes conduct simultaneously; voltage drop
$=2\times0.65=1.30\,\text{V}$.
$V_{peak,out}=17.0-1.30=15.7\,\text{V}$. \quad
(c)~$\tau=RC=1000\times10^{-6}\times100
=0.10\,\text{s}$.
Period of full-wave output $=1/(2f)=1/100
=0.010\,\text{s}$.
Since $\tau\gg T_{output}$ ($0.10\gg0.010$), the
capacitor discharges very little between successive
peaks. Smoothing is highly effective --- the ripple
voltage is small compared with the DC output. \quad
(d)~(i)~Without capacitor: sequence of positive
half-sine pulses at 100\,Hz, peak at $15.7\,\text{V}$,
touching zero at each cycle. Label horizontal axis
``time'', mark peak $V_0=15.7\,\text{V}$. \\
(ii)~With capacitor: nearly flat DC at $\approx15.7\,\text{V}$
with small sawtooth ripple ($\sim0.5$--$1\,\text{V}$
amplitude) at 100\,Hz, showing slow discharge between
charging peaks. \quad
(e)~Full-wave rectification: (1)~both half-cycles
contribute, so average output voltage is approximately
double that of half-wave ($0.637\,V_0$ vs
$0.318\,V_0$); (2)~ripple frequency is $2f=100\,\text{Hz}$
rather than $50\,\text{Hz}$, giving the smoothing
capacitor less time to discharge between peaks and
producing much smaller ripple; (3)~transformer and
circuit components are used more efficiently since
current is drawn during both half-cycles.

\textbf{B3.}
(a)~$I_{C,sat}=(V_{CC}-V_{CE,sat})/R_C
=(9.0-0.2)/1000=8.8\times10^{-3}\,\text{A}
=8.8\,\text{mA}$. \quad
(b)~$I_{B,min}=I_{C,sat}/h_{FE}
=8.8\times10^{-3}/150=58.7\,\mu\text{A}$. \quad
(c)~Voltage across $R_B=V_{in}-V_{BE}=5.0-0.6
=4.4\,\text{V}$.
$R_B=4.4/I_{B,min}=4.4/58.7\times10^{-6}
=75\,000\,\Omega=75\,\text{k}\Omega$
(maximum value to guarantee saturation; use
$\leq75\,\text{k}\Omega$ in practice). \quad
(d)~The relay coil replaces $R_C$. When the input
is HIGH, the collector current $I_C=h_{FE}I_B$
(up to 150 times larger than $I_B$) flows through
the relay coil, energising its electromagnet and
closing (or opening) relay contacts in a separate,
electrically isolated high-power circuit. The
transistor uses a milliamp-level signal to switch
many amperes through a load at a completely different
voltage, with full galvanic isolation between
the control and load circuits. \quad
(e)~In a switch the bias point is either in
\textit{cutoff} ($V_{BE}<0.6\,\text{V}$, $I_C=0$,
transistor OFF) or \textit{saturation}
($V_{CE}\approx0.2\,\text{V}$, $I_C$ limited by
$R_C$ and $V_{CC}$, transistor fully ON). The
operating point is driven hard to either extreme;
there is no linear relationship between input and
output. An amplifier, by contrast, is biased to a
Q-point in the middle of the active (linear) region.
Here $I_C\approx h_{FE}I_B$ throughout the signal
swing, so the output faithfully reproduces and
amplifies the input waveform without clipping.

\textbf{B4.}
(a)~$A_v=-R_f/R_{in}=-47/10=-4.7$. \quad
(b)~(i)~$V_{out}=-4.7\times0.2=-0.94\,\text{V}$. \quad
(ii)~$V_{out}=-4.7\times(-0.3)=+1.41\,\text{V}$. \quad
(iii)~$V_{out}=-4.7\times0.5=-2.35\,\text{V}$.
All three outputs are within $\pm12\,\text{V}$;
the op-amp operates in the linear region throughout.
\quad
(c)~Output saturates when $|A_v\times V_{in}|
\approx12\,\text{V}$:
$V_{in,sat}=12/4.7=2.55\,\text{V}$.
In practice the op-amp saturates slightly below the
rail, so output clips at $\approx\pm(10\text{--}11)\,
\text{V}$ depending on the device; saturation begins
around $|V_{in}|\approx2.3$--$2.6\,\text{V}$. \quad
(d)~Roles swapped: $R_{input}=47\,\text{k}\Omega$,
$R_{feedback}=10\,\text{k}\Omega$.
$A_v=-10/47=-0.213$.
This is an \textit{attenuator} (gain magnitude less
than 1). \quad
(e)~Non-inverting: $A_v=1+R_f/R_1=6$;
$R_f/R_1=5$; $R_1=10\,\text{k}\Omega$;
$R_f=5\times10=\mathbf{50\,\text{k}\Omega}$.

\textbf{B5.}
(a)~$U=(P\cdot K)+(P\cdot T)$. \quad
(b)~Factoring: $U=P(K+T)$
(using Boolean distributive law $AB+AC=A(B+C)$).
This confirms that the PIN correct ($P=1$) is
necessary in all cases; the door unlocks whenever
either $K$ or $T$ (or both) is also HIGH. \quad
(c)~Full truth table:
\begin{center}
\begin{tabular}{ccc|c|c}
$P$ & $K$ & $T$ & $K+T$ & $U=P(K+T)$\\
\hline
0 & 0 & 0 & 0 & 0\\
0 & 0 & 1 & 1 & 0\\
0 & 1 & 0 & 1 & 0\\
0 & 1 & 1 & 1 & 0\\
1 & 0 & 0 & 0 & 0\\
1 & 0 & 1 & 1 & 1\\
1 & 1 & 0 & 1 & 1\\
1 & 1 & 1 & 1 & 1\\
\end{tabular}
\end{center}
The door unlocks only when $P=1$ and at least one of
$K$ or $T$ is HIGH. \quad
(d)~Minimum gate circuit: $K$ and $T$ feed an
\textbf{OR} gate; the OR output and $P$ feed an
\textbf{AND} gate to give $U$. Two gates total. \quad
(e)~Implementing with NAND gates only:
NOT A from NAND: connect both inputs of a NAND
together: $\overline{A}=\text{NAND}(A,A)$.
OR from three NANDs (by De Morgan):
$A+B=\overline{\bar{A}\cdot\bar{B}}$,
so:
$G_1=\text{NAND}(K,K)=\bar{K}$;
$G_2=\text{NAND}(T,T)=\bar{T}$;
$G_3=\text{NAND}(G_1,G_2)=\overline{\bar{K}\cdot\bar{T}}=K+T$.
Then AND from two NANDs:
$G_4=\text{NAND}(P,G_3)=\overline{P\cdot(K+T)}$;
$G_5=\text{NAND}(G_4,G_4)=P(K+T)=U$.
Total: \textbf{five NAND gates}.

\endgroup

\normalsize

\normalsize
